% E1.12 -- a teoria em blocos: o block LASSO balanceado de D44 no desenho de
% produtos. Promove a sondagem de E1.11 (08a-sondagem-blocos.md, inteira, com
% os adendos das §§ 11 e 12) a resultados numerados e provados, para o
% estimador de D44 (níveis c_l livres, os níveis com 2^j < b_n de cada bloco
% num pedaço só, cada nível mais fino em pedaços de b_n a 2 b_n - 1 colunas,
% pesos de razão limitada), seguido de limiar (D45; a cota do limiarizado está
% no adendo de E1.12 de 06-selecao-limiar.tex).
%
% Consome: o viés de E1.3 (02-aproximacao-besov.tex, Lemas 1 e 2 e Corolário
% 1), o autovalor de E1.4 (03-desenho-produtos.tex, Proposição 3), a
% perfilagem e o Lema 6 de E1.5 (04-oraculo.tex) e o Lema 9(iii) e o regime de
% E1.6 (05-taxas.tex). NÃO altera 04 nem 05: os resultados do LASSO ficam para
% a opção coordenada (D43).
%
% Molde: 04-oraculo.tex e 05-taxas.tex (isto é, Bühlmann e van de Geer, 2011,
% Seção 6.2, com a perfilagem de E1.5); a calibração é o Teorema 2.1 de Hsu,
% Kakade e Zhang (2012) com união sobre os pedaços; o risco ideal por pedaços
% tem a forma do Lema 4 do suplemento de Klopp e Pensky (2015).
%
% Numeração global (docs/TAREFA.md, §5): continua de Proposição 6, Lema 13,
% Teorema 2 e Corolário 8, e imprime o número global via \setcounter, como
% 04, 05, 06 e 07. E1.13 acrescentou o Corolário 14 (a taxa lenta em blocos,
% D47(c)), na §5, com o contador posto em 13 antes dele: os Corolários 12 e 13
% estão no adendo de E1.12 de 06-selecao-limiar.tex. E1.14 acrescentou ao
% Corolário 14 os itens (iii) e (iv) (todo s' > 0 e a fronteira s = 1/2, com o
% nível acima do teto e lambda_n^G inflado por (mu^G_J)^{1/2}), sem número novo,
% e a Parte VIII da conferência. E1.15 acrescentou o Corolário 15 (a taxa lenta
% pelo comparador truncado por pedaço, com a contagem de
% sum_G min(lambda w_G ||theta*_G||, ||theta*_G||^2) sob Besov), na §5, depois da
% Observação rem:lenta-desenho, e a Parte IX da conferência.
%
% Conferência numérica: check/08-blocos.R (imprime OK; ~12 min).
% Notação: docs/notacao.md (congelada em E1.1). Os símbolos da teoria em
% blocos (os de E1.11) são macros locais deste arquivo, como em 07, até
% entrarem no notacao.md (§10) pelo chat principal.
\documentclass[11pt,a4paper]{article}
\usepackage[T1]{fontenc}
\usepackage[utf8]{inputenc}
\usepackage[brazil]{babel}
\usepackage[margin=2.6cm]{geometry}
\usepackage{enumitem}
\usepackage{booktabs}
\usepackage{xcolor}
\usepackage[colorlinks=true,linkcolor=blue!60!black,citecolor=blue!60!black]{hyperref}
% Macros do WAFC. Implementa docs/notacao.md (congelada em E1.1, 2026-09-18).
% Único lugar onde uma macro é definida; os derivations/*.tex fazem % Macros do WAFC. Implementa docs/notacao.md (congelada em E1.1, 2026-09-18).
% Único lugar onde uma macro é definida; os derivations/*.tex fazem % Macros do WAFC. Implementa docs/notacao.md (congelada em E1.1, 2026-09-18).
% Único lugar onde uma macro é definida; os derivations/*.tex fazem \input{macros}.
\usepackage{amsmath,amssymb,amsthm}

% --- básicas -------------------------------------------------------------
\newcommand{\R}{\mathbb{R}}
% D24: a §5 das instruções da Statistica Sinica manda esperança e
% probabilidade em romano; o manuscrito já nasceu assim. Os delimitadores
% ficam como cada arquivo os escreveu (colchetes em vários lugares): o que
% D24 exige e que aqui se cumpre é o símbolo, não a troca de todos os
% delimitadores em texto de trabalho.
\newcommand{\E}{\mathrm{E}}
\newcommand{\Prob}{\mathrm{P}}
\newcommand{\Var}{\mathrm{Var}}
\newcommand{\T}{^{\top}}
\newcommand{\norm}[1]{\left\lVert #1 \right\rVert}
\newcommand{\normn}[1]{\left\lVert #1 \right\rVert_{n}}
\newcommand{\normL}[1]{\left\lVert #1 \right\rVert_{L_2(P)}}
\newcommand{\inner}[2]{\langle #1, #2 \rangle}

% --- índices (notacao.md, §1) --------------------------------------------
% A covariável linear é X_\ell, \ell = 1..p; a moduladora é U_m, m = 1..q;
% a wavelet é \psi_{jk}, com nível j e translação k. O par (\ell,m) é o bloco.
\newcommand{\cv}{\ell}          % índice da covariável linear
\newcommand{\md}{m}             % índice da covariável moduladora
\newcommand{\lev}{j}            % nível de resolução
\newcommand{\tsl}{k}            % translação
\newcommand{\jzero}{j_{0}}
\newcommand{\NJ}{N_{J}}

% --- modelo (notacao.md, §2) ---------------------------------------------
\newcommand{\bX}{\mathbf{X}}
\newcommand{\bU}{\mathbf{U}}
\newcommand{\bZ}{\mathbf{Z}}
\newcommand{\fcoef}[1]{\beta_{#1}}          % \fcoef{\cv}: coeficiente funcional
\newcommand{\lvl}[1]{c_{#1}}                % \lvl{\cv}: nível não penalizado
\newcommand{\gcomp}[2]{g_{#1#2}}            % \gcomp{\cv}{\md}: componente aditiva

% --- base e espaço de aproximação (notacao.md, §3) -----------------------
\newcommand{\wav}[2]{\psi_{#1#2}}           % \wav{\lev}{\tsl}
\newcommand{\VJ}{V_{J}}
\newcommand{\WJ}{W_{J}}
\newcommand{\PiJ}{\Pi_{J}}
\newcommand{\coef}[4]{\theta_{#1#2,#3#4}}   % \coef{\cv}{\md}{\lev}{\tsl}
\newcommand{\coefs}[4]{\theta^{*}_{#1#2,#3#4}}
\newcommand{\blk}[2]{\theta_{#1#2,\cdot}}   % subvetor do bloco (grupo em E2.2)
\newcommand{\btheta}{\boldsymbol{\theta}}
\newcommand{\thetastar}{\btheta^{*}}
\newcommand{\thetahat}{\widehat{\btheta}}
\newcommand{\bc}{\mathbf{c}}
\newcommand{\chat}{\widehat{\mathbf{c}}}
\newcommand{\Gram}{\Sigma}
\newcommand{\Gramhat}{\widehat{\Sigma}}

% --- estimador (notacao.md, §4) ------------------------------------------
\newcommand{\pen}{\lambda}
\newcommand{\penn}{\lambda_{n}}
\newcommand{\wgt}[4]{w_{#1#2,#3#4}}
\newcommand{\compat}{\phi_{0}}
\newcommand{\Szero}{S_{0}}
% A penalidade é escrita \norm{\btheta}_{1}; nunca "\ell_1", que colidiria
% com o índice \cv (notacao.md, §1).

% --- espaços de funções (notacao.md, §5) ---------------------------------
\newcommand{\Besov}[3]{B^{#1}_{#2,#3}}
\newcommand{\seff}{s'}
\newcommand{\weakl}[1]{\mathrm{w}\ell_{#1}}

% --- ambientes -----------------------------------------------------------
% D29: derivações em português, código em inglês. Os nomes dos ambientes
% seguem a prosa dos arquivos; o manuscrito não usa este arquivo (as macros
% estão copiadas para dentro dele, e o template da revista define os seus
% próprios ambientes, em inglês).
\newtheorem{proposition}{Proposição}
\newtheorem{lemma}{Lema}
\newtheorem{theorem}{Teorema}
\newtheorem{corollary}{Corolário}
\theoremstyle{definition}
\newtheorem{assumption}{Hipótese}
\newtheorem{remark}{Observação}
.
\usepackage{amsmath,amssymb,amsthm}

% --- básicas -------------------------------------------------------------
\newcommand{\R}{\mathbb{R}}
% D24: a §5 das instruções da Statistica Sinica manda esperança e
% probabilidade em romano; o manuscrito já nasceu assim. Os delimitadores
% ficam como cada arquivo os escreveu (colchetes em vários lugares): o que
% D24 exige e que aqui se cumpre é o símbolo, não a troca de todos os
% delimitadores em texto de trabalho.
\newcommand{\E}{\mathrm{E}}
\newcommand{\Prob}{\mathrm{P}}
\newcommand{\Var}{\mathrm{Var}}
\newcommand{\T}{^{\top}}
\newcommand{\norm}[1]{\left\lVert #1 \right\rVert}
\newcommand{\normn}[1]{\left\lVert #1 \right\rVert_{n}}
\newcommand{\normL}[1]{\left\lVert #1 \right\rVert_{L_2(P)}}
\newcommand{\inner}[2]{\langle #1, #2 \rangle}

% --- índices (notacao.md, §1) --------------------------------------------
% A covariável linear é X_\ell, \ell = 1..p; a moduladora é U_m, m = 1..q;
% a wavelet é \psi_{jk}, com nível j e translação k. O par (\ell,m) é o bloco.
\newcommand{\cv}{\ell}          % índice da covariável linear
\newcommand{\md}{m}             % índice da covariável moduladora
\newcommand{\lev}{j}            % nível de resolução
\newcommand{\tsl}{k}            % translação
\newcommand{\jzero}{j_{0}}
\newcommand{\NJ}{N_{J}}

% --- modelo (notacao.md, §2) ---------------------------------------------
\newcommand{\bX}{\mathbf{X}}
\newcommand{\bU}{\mathbf{U}}
\newcommand{\bZ}{\mathbf{Z}}
\newcommand{\fcoef}[1]{\beta_{#1}}          % \fcoef{\cv}: coeficiente funcional
\newcommand{\lvl}[1]{c_{#1}}                % \lvl{\cv}: nível não penalizado
\newcommand{\gcomp}[2]{g_{#1#2}}            % \gcomp{\cv}{\md}: componente aditiva

% --- base e espaço de aproximação (notacao.md, §3) -----------------------
\newcommand{\wav}[2]{\psi_{#1#2}}           % \wav{\lev}{\tsl}
\newcommand{\VJ}{V_{J}}
\newcommand{\WJ}{W_{J}}
\newcommand{\PiJ}{\Pi_{J}}
\newcommand{\coef}[4]{\theta_{#1#2,#3#4}}   % \coef{\cv}{\md}{\lev}{\tsl}
\newcommand{\coefs}[4]{\theta^{*}_{#1#2,#3#4}}
\newcommand{\blk}[2]{\theta_{#1#2,\cdot}}   % subvetor do bloco (grupo em E2.2)
\newcommand{\btheta}{\boldsymbol{\theta}}
\newcommand{\thetastar}{\btheta^{*}}
\newcommand{\thetahat}{\widehat{\btheta}}
\newcommand{\bc}{\mathbf{c}}
\newcommand{\chat}{\widehat{\mathbf{c}}}
\newcommand{\Gram}{\Sigma}
\newcommand{\Gramhat}{\widehat{\Sigma}}

% --- estimador (notacao.md, §4) ------------------------------------------
\newcommand{\pen}{\lambda}
\newcommand{\penn}{\lambda_{n}}
\newcommand{\wgt}[4]{w_{#1#2,#3#4}}
\newcommand{\compat}{\phi_{0}}
\newcommand{\Szero}{S_{0}}
% A penalidade é escrita \norm{\btheta}_{1}; nunca "\ell_1", que colidiria
% com o índice \cv (notacao.md, §1).

% --- espaços de funções (notacao.md, §5) ---------------------------------
\newcommand{\Besov}[3]{B^{#1}_{#2,#3}}
\newcommand{\seff}{s'}
\newcommand{\weakl}[1]{\mathrm{w}\ell_{#1}}

% --- ambientes -----------------------------------------------------------
% D29: derivações em português, código em inglês. Os nomes dos ambientes
% seguem a prosa dos arquivos; o manuscrito não usa este arquivo (as macros
% estão copiadas para dentro dele, e o template da revista define os seus
% próprios ambientes, em inglês).
\newtheorem{proposition}{Proposição}
\newtheorem{lemma}{Lema}
\newtheorem{theorem}{Teorema}
\newtheorem{corollary}{Corolário}
\theoremstyle{definition}
\newtheorem{assumption}{Hipótese}
\newtheorem{remark}{Observação}
.
\usepackage{amsmath,amssymb,amsthm}

% --- básicas -------------------------------------------------------------
\newcommand{\R}{\mathbb{R}}
% D24: a §5 das instruções da Statistica Sinica manda esperança e
% probabilidade em romano; o manuscrito já nasceu assim. Os delimitadores
% ficam como cada arquivo os escreveu (colchetes em vários lugares): o que
% D24 exige e que aqui se cumpre é o símbolo, não a troca de todos os
% delimitadores em texto de trabalho.
\newcommand{\E}{\mathrm{E}}
\newcommand{\Prob}{\mathrm{P}}
\newcommand{\Var}{\mathrm{Var}}
\newcommand{\T}{^{\top}}
\newcommand{\norm}[1]{\left\lVert #1 \right\rVert}
\newcommand{\normn}[1]{\left\lVert #1 \right\rVert_{n}}
\newcommand{\normL}[1]{\left\lVert #1 \right\rVert_{L_2(P)}}
\newcommand{\inner}[2]{\langle #1, #2 \rangle}

% --- índices (notacao.md, §1) --------------------------------------------
% A covariável linear é X_\ell, \ell = 1..p; a moduladora é U_m, m = 1..q;
% a wavelet é \psi_{jk}, com nível j e translação k. O par (\ell,m) é o bloco.
\newcommand{\cv}{\ell}          % índice da covariável linear
\newcommand{\md}{m}             % índice da covariável moduladora
\newcommand{\lev}{j}            % nível de resolução
\newcommand{\tsl}{k}            % translação
\newcommand{\jzero}{j_{0}}
\newcommand{\NJ}{N_{J}}

% --- modelo (notacao.md, §2) ---------------------------------------------
\newcommand{\bX}{\mathbf{X}}
\newcommand{\bU}{\mathbf{U}}
\newcommand{\bZ}{\mathbf{Z}}
\newcommand{\fcoef}[1]{\beta_{#1}}          % \fcoef{\cv}: coeficiente funcional
\newcommand{\lvl}[1]{c_{#1}}                % \lvl{\cv}: nível não penalizado
\newcommand{\gcomp}[2]{g_{#1#2}}            % \gcomp{\cv}{\md}: componente aditiva

% --- base e espaço de aproximação (notacao.md, §3) -----------------------
\newcommand{\wav}[2]{\psi_{#1#2}}           % \wav{\lev}{\tsl}
\newcommand{\VJ}{V_{J}}
\newcommand{\WJ}{W_{J}}
\newcommand{\PiJ}{\Pi_{J}}
\newcommand{\coef}[4]{\theta_{#1#2,#3#4}}   % \coef{\cv}{\md}{\lev}{\tsl}
\newcommand{\coefs}[4]{\theta^{*}_{#1#2,#3#4}}
\newcommand{\blk}[2]{\theta_{#1#2,\cdot}}   % subvetor do bloco (grupo em E2.2)
\newcommand{\btheta}{\boldsymbol{\theta}}
\newcommand{\thetastar}{\btheta^{*}}
\newcommand{\thetahat}{\widehat{\btheta}}
\newcommand{\bc}{\mathbf{c}}
\newcommand{\chat}{\widehat{\mathbf{c}}}
\newcommand{\Gram}{\Sigma}
\newcommand{\Gramhat}{\widehat{\Sigma}}

% --- estimador (notacao.md, §4) ------------------------------------------
\newcommand{\pen}{\lambda}
\newcommand{\penn}{\lambda_{n}}
\newcommand{\wgt}[4]{w_{#1#2,#3#4}}
\newcommand{\compat}{\phi_{0}}
\newcommand{\Szero}{S_{0}}
% A penalidade é escrita \norm{\btheta}_{1}; nunca "\ell_1", que colidiria
% com o índice \cv (notacao.md, §1).

% --- espaços de funções (notacao.md, §5) ---------------------------------
\newcommand{\Besov}[3]{B^{#1}_{#2,#3}}
\newcommand{\seff}{s'}
\newcommand{\weakl}[1]{\mathrm{w}\ell_{#1}}

% --- ambientes -----------------------------------------------------------
% D29: derivações em português, código em inglês. Os nomes dos ambientes
% seguem a prosa dos arquivos; o manuscrito não usa este arquivo (as macros
% estão copiadas para dentro dele, e o template da revista define os seus
% próprios ambientes, em inglês).
\newtheorem{proposition}{Proposição}
\newtheorem{lemma}{Lema}
\newtheorem{theorem}{Teorema}
\newtheorem{corollary}{Corolário}
\theoremstyle{definition}
\newtheorem{assumption}{Hipótese}
\newtheorem{remark}{Observação}


% --- atalhos locais deste arquivo (como em 04, 05 e 07) --------------------
\newcommand{\lmin}{\lambda_{\min}}
\newcommand{\lmaxx}{\lambda_{\max}}
\newcommand{\Amat}{\mathbf{A}}
\newcommand{\Bmat}{\mathbf{B}}
\newcommand{\Btil}{\widetilde{\Bmat}}
\newcommand{\PA}{\Pi_{\Amat}}
\newcommand{\MA}{M_{\Amat}}
\newcommand{\bbeta}{\boldsymbol{\beta}}
\newcommand{\betahat}{\widehat{\bbeta}}
\newcommand{\betastar}{\bbeta^{*}}
\newcommand{\bveps}{\boldsymbol{\varepsilon}}
\newcommand{\bb}{\mathbf{b}}
\newcommand{\bY}{\mathbf{Y}}
\newcommand{\smax}{\widehat{\sigma}_{\max}}
\newcommand{\gtil}{\widetilde{\gamma}}
\newcommand{\GBA}{\Gramhat_{\Bmat\mid\Amat}}
\newcommand{\thbar}{\bar{\btheta}}
\newcommand{\bbar}{\bar{\bb}}
\newcommand{\vbar}{\bar{v}}
\newcommand{\opnorm}[1]{\left\lVert #1 \right\rVert_{\mathrm{op}}}
% --- símbolos da teoria em blocos (E1.11; proposta para o notacao.md, §10) --
\newcommand{\Gcal}{\mathcal{G}}
\newcommand{\Gzero}{\mathcal{G}_{0}}
\newcommand{\Gbar}{\bar{\mathcal{G}}}
\newcommand{\Hcal}{\mathcal{H}}
\newcommand{\bn}{b_{n}}
\newcommand{\nGw}[1]{\lVert #1\rVert_{\mathcal{G},w}}
\newcommand{\nGone}[1]{\lVert #1\rVert_{\mathcal{G},1}}
\newcommand{\Psitil}{\widetilde{\Psi}}
\newcommand{\Rid}{R_{\mathcal{G}}}
\newcommand{\lamG}{\lambda^{\mathcal{G}}_{n}}
\newcommand{\lamGone}{\lambda^{\mathcal{G}}_{n,1}}
\newcommand{\lamGp}{\lambda^{\mathcal{G},+}_{n}}   % E1.14: \lamG inflado acima do teto
\newcommand{\lamz}[1]{\lambda_{0,#1}}
\newcommand{\lamzbar}[1]{\bar{\lambda}_{0,#1}}
\newcommand{\rhoG}{\rho^{\mathcal{G}}_{n}}
\newcommand{\rhoGt}{\tilde{\rho}^{\mathcal{G}}_{n}}
\newcommand{\Tw}{\mathcal{T}_{\mathcal{G},w}}
\newcommand{\Ecal}{\mathcal{E}}
\newcommand{\ESig}{\mathcal{E}_{\Sigma}}
\newcommand{\Acal}{\mathcal{A}_{n}}
\newcommand{\jst}{j^{*}}
\newcommand{\AG}{A^{\mathcal{G}}_{s,\pi}}

% Numeração global (docs/TAREFA.md, §5): Proposição 7, Lema 14, Teorema 3 e
% Corolário 9 são os próximos números livres; os contadores abaixo fazem o
% número impresso JÁ SER o global.
\setcounter{proposition}{6}
\setcounter{lemma}{13}
\setcounter{theorem}{2}
\setcounter{corollary}{8}

% A hipótese de regime deste arquivo ganha a letra B (blocos), como a
% Hipótese S de 06, para não colidir com a Hipótese 1 de E1.6, citada ao lado.
\renewcommand{\theassumption}{B}

\title{E1.12. A teoria em blocos: o block LASSO balanceado no desenho de
produtos}
\author{wafc-draft, \texttt{derivations/08-blocos.tex}}
\date{2026-10-03}

\begin{document}
\maketitle

\section*{Onde este resultado entra}

A decisão D44 fez do WAFC o block LASSO na forma balanceada, seguido de limiar
nas normas por bloco (D45), e deixou o LASSO coordenado de D3 como opção
(D43). A teoria de E1.5 e E1.6 (\texttt{04-oraculo.tex} e
\texttt{05-taxas.tex}) é do LASSO. A sondagem de E1.11
(\texttt{08a-sondagem-blocos.md}) mostrou que ela transfere para a penalidade
em blocos sem cone e com as mesmas constantes: muda a calibração, a
compressibilidade dá lugar ao risco ideal por pedaços, e a taxa ganha um
fator logarítmico; os adendos das §§~11 e~12 mostraram que pesos de razão
limitada custam só uma constante, o que cobre os pesos $\sqrt{|G|}$ do
\texttt{grpreg} na forma balanceada. Este arquivo escreve isso como
resultados numerados e provados. São seis peças:

\begin{enumerate}[label=(\alph*), leftmargin=*, nosep]
  \item a \textbf{forma balanceada} (Proposição~\ref{prop:balanceada}): os
    tamanhos dos pedaços, o número deles e a razão $\varrho^{2} \le 3$ dos pesos
    $\sqrt{|G|}$;
  \item a \textbf{calibração em blocos} (Lema~\ref{lem:calib-blocos}): Hsu,
    Kakade e Zhang (2012) com união sobre os pedaços, um evento que não
    depende dos pesos, e o preço $\varrho$ dos pesos;
  \item o \textbf{oráculo sem cone} (Teorema~\ref{thm:oraculo-blocos}),
    escrito por inteiro e com comparador arbitrário, e a sua \textbf{forma de
    risco ideal} (Corolário~\ref{cor:risco-ideal});
  \item a \textbf{aproximação}: a hipótese de Besov de E1.3 implica a cota do
    risco ideal por pedaços, sem hipótese nova, nos dois regimes de $\pi$
    (Lema~\ref{lem:risco-pedacos});
  \item as \textbf{taxas}: a do espaço de aproximação sem logaritmo
    (Teorema~\ref{thm:taxa-blocos}), a das componentes
    (Corolário~\ref{cor:componentes-blocos}) e a adaptativa
    $n^{-2s/(2s+1)}(\log n)^{(2/\pi-1)_{+}/(2s+1)}$, sem logaritmo em
    $\pi \ge 2$ (Corolário~\ref{cor:adaptativa-blocos});
  \item a \textbf{taxa lenta}, sem condição de desenho
    (Corolário~\ref{cor:lenta-blocos}, acrescentado em E1.13), a versão em
    blocos da Proposição~4 de E1.6: a mesma ordem pela cadeia com
    $\norm{\thetastar}_{1}$, e sem o logaritmo em $s < 1/2$ e $\pi \ge 2$
    pela contagem na norma de grupos.
\end{enumerate}

E1.15 acrescentou (g), a \textbf{taxa lenta pelo comparador truncado por
pedaço} (Corolário~\ref{cor:truncada-blocos}): em $s < 1/2$ ela alcança a taxa
do Corolário~\ref{cor:adaptativa-blocos} sem condição de desenho, e em
$s \ge 1/2$ não muda a de (f).

Pelo D16 o enunciado principal do artigo é o Corolário~5 de E1.6. Com D44 o
candidato passa a ser o Corolário~\ref{cor:adaptativa-blocos}, que é o mesmo
enunciado com a penalidade em blocos e um fator $(\log n)^{2\seff/(2s+1)}$ a
menos; a confirmação é do autor, na integração. O Corolário~8 com o ajuste
em blocos e a cota de risco do estimador limiarizado estão no adendo de E1.12
de \texttt{06-selecao-limiar.tex} (Lema~16 e Corolários~12 e~13).

\section{Objetos}

\paragraph{Modelo, amostra e desenho.} Como em E1.5, \texttt{04-oraculo.tex},
§1: $(Y_{i},\bX_{i},\bU_{i})$ i.i.d.\ com $Y_{i} = f(\bX_{i},\bU_{i}) +
\varepsilon_{i}$, $f(x,u) = \sum_{\cv}x_{\cv}\{\lvl{\cv} +
\sum_{\md}\gcomp{\cv}{\md}(u_{\md})\}$; na ordem de colunas de D12,
$\bZ = [\Amat\;\Bmat]$ com $\Amat_{i\cv} = X_{i\cv}$ e
$\Bmat_{i,(\cv\md,\lev\tsl)} = X_{i\cv}\wav{\lev}{\tsl}(U_{i\md})$,
$d = pq\NJ$, $\NJ = 2^{J}-1$, e $\Gramhat = \bZ\T\bZ/n$.

\paragraph{Pedaços e pesos.} $\Gcal$ é uma partição do conjunto das $d$
coordenadas penalizadas em \emph{pedaços} $G$, cada um contido num único bloco
$(\cv,\md)$; $|\Gcal|$ é o número de pedaços, $|G|$ o de colunas do pedaço, e
$\btheta_{G} \in \R^{|G|}$ o subvetor. Os pesos são $w = (w_{G})_{G\in\Gcal}$,
$w_{G} > 0$, com razão
\[
  \varrho \;=\; \frac{\max_{G}w_{G}}{\min_{G}w_{G}} \;\ge\; 1 ,
\]
e a penalidade é a norma
\[
  \nGw{\btheta} \;=\; \sum_{G\in\Gcal} w_{G}\,\norm{\btheta_{G}}_{2},
  \qquad
  \nGw{\btheta_{\Hcal}} \;=\; \sum_{G\in\Hcal} w_{G}\,\norm{\btheta_{G}}_{2}
  \quad (\Hcal \subseteq \Gcal).
\]
Para $\Hcal \subseteq \Gcal$ escrevemos $\mathcal{W}(\Hcal) = \sum_{G\in\Hcal}w_{G}^{2}$,
e para $\btheta \in \R^{d}$, $\Gzero(\btheta) = \{G : \btheta_{G} \ne 0\}$, os
pedaços ativos. Com $w_{G} = 1$, $\mathcal{W}(\Gzero(\btheta))$ é o número de pedaços
ativos; com $w_{G} = \sqrt{|G|}$, é o número de coordenadas nos pedaços
ativos.

\paragraph{Estimador.} Com $\pen > 0$,
\begin{equation}
  \label{eq:est-blocos}
  \betahat = (\chat,\thetahat)
  \;\in\; \operatorname*{arg\,min}_{(\bc,\btheta)\in\R^{p}\times\R^{d}}
  \Bigl\{\normn{\bY - \Amat\bc - \Bmat\btheta}^{2} + 2\pen\nGw{\btheta}\Bigr\},
  \qquad \widehat f = \Amat\chat + \Bmat\thetahat .
\end{equation}
Os níveis $\lvl{\cv}$ não são penalizados (D3). Na normalização de Klopp e
Pensky (2015, eq.~3.3), $n^{-1}\norm{\bY - \Bmat\alpha}_{2}^{2} +
\delta\norm{\alpha}_{\mathrm{block}}$, vale $\delta = 2\pen$; na do
\texttt{grpreg}, $(2n)^{-1}\norm{\cdot}_{2}^{2} + \pen\cdot$penalidade, o $\pen$
é o mesmo.

\paragraph{A forma balanceada (D44).} Fixe um inteiro $b \ge 2$ e seja
$\jst = \max\{\lev \ge 0 : 2^{\lev} < b\}$. Em cada bloco $(\cv,\md)$:
\begin{enumerate}[label=(\roman*), leftmargin=*, nosep]
  \item os níveis $\lev \le \min(\jst, J-1)$ formam um único pedaço, o
    \emph{pedaço grosso};
  \item cada nível $\lev$ com $\jst < \lev \le J-1$ (os \emph{níveis finos},
    $2^{\lev} \ge b$) é cortado em $k_{\lev} = \lfloor 2^{\lev}/b\rfloor$
    pedaços de translações consecutivas, os $k_{\lev}-1$ primeiros com $b$
    colunas e o último com as $b + (2^{\lev} - k_{\lev}b)$ restantes.
\end{enumerate}
O WAFC de D44 é \eqref{eq:est-blocos} com essa partição, $b = \bn =
\lceil\log n\rceil$ e $w_{G} = \sqrt{|G|}$; no código, é o
\texttt{wafc\_kp\_groups(balanced = TRUE)} com o \texttt{group.multiplier}
padrão do \texttt{grpreg}, que resolve a variante branca da
Observação~\ref{rem:branca}.

\paragraph{Perfilagem.} Como em E1.5: $\PA$ é a projeção sobre o espaço-coluna
de $\Amat$, $\MA = I_{n} - \PA$, $\Btil = \MA\Bmat$,
$\widetilde\bY = \MA\bY$, $\GBA = \Btil\T\Btil/n$ e
$\gtil = \lmin(\GBA)$; por pedaço,
\[
  \Psitil_{G} \;=\; \Btil_{G}\T\Btil_{G}/n \in \R^{|G|\times|G|},
\]
com $\Btil_{G}$ as colunas de $\Btil$ em $G$; e
$\smax^{2} = \max_{a}(\Gramhat)_{aa}$ sobre as colunas penalizadas.

\paragraph{Constantes e eventos de E1.4 e E1.6.} $\gamma = \kappa_{1}c_{U}$ e
$\Lambda = \kappa_{2}C_{U} + \gamma/2$, como em E1.6, e
\[
  \ESig \;=\; \bigl\{\opnorm{\Gramhat - \Gram} < \gamma/2\bigr\}.
\]
No evento $\ESig$, a Proposição~3(iii) de E1.4 e Weyl dão
$\lmin(\Gramhat) > \gamma/2$ e $\lmaxx(\Gramhat) < \Lambda$ (é o Lema~9(iii) de
E1.6), e a Proposição~3(iv) dá $\Prob(\ESig) \to 1$ quando
$pq2^{J}\log(pq2^{J})/n \to 0$.

\paragraph{Oráculo, viés e comparadores.} $\betastar = (\bc,\thetastar)$ é o
oráculo de E1.3, com $\coefs{\cv}{\md}{\lev}{\tsl} =
\inner{\gcomp{\cv}{\md}}{\wav{\lev}{\tsl}}$ para $\lev < J$, $f_{J} =
\bZ\betastar$, $\bb = f - f_{J}$, e
$\mathcal{B}_{J}^{2} = (pq)^{2}B_{X}^{2}C_{U}C_{g}^{2}(1-2^{-2\seff})^{-1}
2^{-2J\seff}$ é a cota de $\E\normn{\bb}^{2}$ do Corolário~1 de E1.3 (com o
$C_{X}$ de lá escrito $B_{X}$, D24). Um \emph{comparador} é qualquer
$\bar\bbeta = (\bar\bc,\thbar) \in \R^{p}\times\R^{d}$, com resíduo
determinístico $\bbar = f - \bZ\bar\bbeta$.

\paragraph{Risco ideal por pedaços.} Para $\btheta \in \R^{d}$ e $\eta > 0$,
\[
  \Rid(\btheta;\eta) \;=\; \sum_{G\in\Gcal}\min\bigl(\norm{\btheta_{G}}_{2}^{2},\,\eta\bigr).
\]
É não decrescente em $\eta$, e como $\min(u,\kappa\eta) \le \kappa\min(u,\eta)$
para $\kappa \ge 1$, vale $\Rid(\btheta;\kappa\eta) \le \kappa\Rid(\btheta;\eta)$.
Com $\Gcal$ de pedaços unitários, $\Rid$ é o risco ideal da projeção
coordenada a coordenada.

\section{Hipóteses e o que se usa de E1.5}

Além das hipóteses de E1.3 (base, Besov, densidade das moduladoras, $\bX$
limitada), de E1.4 (moduladoras, covariáveis lineares, base) e das
Hipóteses~1 (amostra i.i.d.) e~2 (erro sub-gaussiano de parâmetro $\sigma$) de
E1.5, citadas onde entram, este arquivo usa uma hipótese de regime, a
Hipótese~\ref{ass:regime-blocos} da §\ref{sec:taxas}, só nos enunciados
assintóticos.

O Lema~4 de E1.5 é usado com a penalidade em blocos. Ele não consome número
novo: o contador abaixo é recuado para que apareça com o número que já tem.

\setcounter{lemma}{3}
\begin{lemma}[perfilagem; provado em E1.5 para $\norm{\cdot}_{1}$]\label{lem:perfil-blocos}
Suponha $\Amat$ de posto $p$. Com a penalidade $\nGw{\cdot}$ em
\eqref{eq:est-blocos}:
\begin{enumerate}[label=\textup{(\roman*)}, leftmargin=*, nosep]
  \item $\thetahat$ minimiza $\normn{\widetilde\bY - \Btil\btheta}^{2} +
    2\pen\nGw{\btheta}$, e $\chat = (\Amat\T\Amat)^{-1}\Amat\T(\bY -
    \Bmat\thetahat)$;
  \item para todo comparador $\bar\bbeta$, com $\vbar = \thetahat - \thbar$,
    \begin{equation}
      \label{eq:decomp-blocos}
      \widehat f - \bZ\bar\bbeta \;=\; \Btil\vbar + \PA(\bbar + \bveps);
    \end{equation}
  \item $\gtil \ge \lmin(\Gramhat)$;
  \item se $\lmin(\Gramhat) > 0$, o minimizador em \eqref{eq:est-blocos} é
    único.
\end{enumerate}
\end{lemma}
\setcounter{lemma}{13}

\begin{proof}
Em E1.5, a prova de (i) usa só que a penalidade não depende de $\bc$
(minimizar primeiro em $\bc$ dá $\normn{\MA(\bY - \Bmat\btheta)}^{2}$), e a de
(iv) só que ela é convexa; $\nGw{\cdot}$ tem as duas propriedades. (iii) não
envolve a penalidade. Para (ii): por (i), $\Amat\chat = \PA(\bY -
\Bmat\thetahat)$, e $\bY = \Amat\bar\bc + \Bmat\thbar + \bbar + \bveps$ pela
definição de $\bbar$; logo $\Amat(\chat - \bar\bc) = \PA(\Bmat\thbar + \bbar +
\bveps - \Bmat\thetahat) = \PA(\bbar + \bveps) - \PA\Bmat\vbar$, e
$\widehat f - \bZ\bar\bbeta = \Amat(\chat - \bar\bc) + \Bmat\vbar =
\MA\Bmat\vbar + \PA(\bbar + \bveps)$.
\end{proof}

\section{A forma balanceada e a calibração}

\begin{proposition}[a forma balanceada]\label{prop:balanceada}
Seja $\Gcal$ a forma balanceada com $b \ge 2$ e $J \ge 1$. Então:
\begin{enumerate}[label=\textup{(\roman*)}, leftmargin=*, nosep]
  \item todo pedaço fino está contido num nível e tem entre $b$ e $2b-1$
    colunas; se $J > \jst + 1$, o pedaço grosso tem $2^{\jst+1}-1$ colunas, e
    $b - 1 \le 2^{\jst+1}-1 \le 2b-3$; se $J \le \jst + 1$, cada bloco é um
    pedaço só, de $2^{J}-1$ colunas;
  \item $|\Gcal| \le pq\,(1 + 2^{J}/b)$;
  \item com $w_{G} = \sqrt{|G|}$, $\varrho^{2} = \max_{G}|G|/\min_{G}|G| \le
    (2b-1)/(b-1) \le 3$; em $b = 6$ e $b = 7$ (os $\bn$ de $n = 250$ e de
    $n = 500$ e $1000$), $\varrho^{2} \le 5/3$ e $\varrho^{2} \le 11/7$ quando
    $J \ge 5$.
\end{enumerate}
\end{proposition}

\begin{lemma}[calibração em blocos]\label{lem:calib-blocos}
Sob as Hipóteses~1 e~2 de E1.5, seja $\Gcal$ uma partição qualquer das
coordenadas penalizadas, e $\alpha \in (0,1)$.
\begin{enumerate}[label=\textup{(\roman*)}, leftmargin=*]
  \item Com
    \[
      \lamz{G} \;=\; \frac{\sigma}{\sqrt n}\Bigl(\sqrt{\operatorname{tr}\Psitil_{G}}
      + \sqrt{2\opnorm{\Psitil_{G}}\log(|\Gcal|/\alpha)}\Bigr),
    \]
    o evento $\Ecal = \bigl\{\norm{(\Btil\T\bveps/n)_{G}}_{2} \le \lamz{G}
    \text{ para todo } G\bigr\}$ tem $\Prob(\Ecal \mid \bX,\bU) \ge 1 -
    \alpha$. Em consequência, para quaisquer pesos $w$ e todo
    $\pen \ge \pen_{w} := 2\max_{G}\lamz{G}/w_{G}$,
    \[
      \Ecal \;\subseteq\; \Tw \;=\; \Bigl\{\max_{G}\,
      \norm{(\Btil\T\bveps/n)_{G}}_{2}/w_{G} \le \pen/2\Bigr\},
      \qquad \Prob(\Tw) \ge 1 - \alpha .
    \]
  \item \textup{(forma fechada)} $\operatorname{tr}\Psitil_{G} \le
    |G|\smax^{2}$ e $\opnorm{\Psitil_{G}} \le \lmaxx(\Gramhat)$; logo, no
    evento $\{\smax^{2} \le 2B_{X}^{2}C_{U}\} \cap \{\lmaxx(\Gramhat) \le
    \Lambda\}$, vale $\lamz{G} \le \lamzbar{G}$ para todo $G$, com
    \[
      \lamzbar{G} \;=\; \frac{\sigma}{\sqrt n}\Bigl(B_{X}\sqrt{2C_{U}|G|}
      + \sqrt{2\Lambda\log(|\Gcal|/\alpha)}\Bigr),
    \]
    que é determinístico.
  \item \textup{(o preço dos pesos)} Com
    \[
      \lamG \;=\; 2\max_{G}\lamzbar{G}/w_{G},
      \qquad
      \lamGone \;=\; 2\max_{G}\lamzbar{G},
    \]
    vale $\lamG w_{G} \le \varrho\,\lamGone$ para todo $G$, e portanto
    $(\lamG)^{2}\mathcal{W}(\Hcal) \le \varrho^{2}(\lamGone)^{2}|\Hcal|$ para todo
    $\Hcal \subseteq \Gcal$. O mesmo vale com $\pen_{w}$ e
    $2\max_{G}\lamz{G}$ no lugar de $\lamG$ e $\lamGone$.
\end{enumerate}
\end{lemma}

O evento $\Ecal$ não sabe dos pesos: o peso só escolhe qual pedaço dita o
$\pen$, o de maior $\lamz{G}/w_{G}$, e o preço de pesos desiguais é o fator
$\varrho$ de (iii), pago no termo de estimação. Com $w_{G} = 1$ e pedaços de até
$b_{\max}$ colunas, $\lamGone \asymp \sigma\sqrt{(b_{\max} +
\log|\Gcal|)/n}$, que é o $\widehat\delta$ de Klopp e Pensky (2015, eq.~3.14) a
menos de constantes (com $p$ fixo, as duas escalas são
$\sigma\sqrt{\log n/n}$; §\ref{sec:kp}) e a versão sub-gaussiana do
$\lambda_{j}$ do Lema~3.1 de Lounici et al.\ (2011).

\section{O oráculo sem cone}

\begin{theorem}[desigualdade oráculo em blocos]\label{thm:oraculo-blocos}
Sob as Hipóteses~1 e~2 de E1.5, seja $\Gcal$ uma partição das coordenadas
penalizadas em pedaços contidos em blocos, $w$ pesos positivos, $\Amat$ de
posto $p$ e $\pen > 0$. No evento $\Tw$ do Lema~\ref{lem:calib-blocos}, que
tem probabilidade ao menos $1-\alpha$ sempre que $\pen \ge \pen_{w}$, vale
\emph{simultaneamente para todo comparador} $\bar\bbeta = (\bar\bc,\thbar)$,
com
\[
  \bbar = f - \bZ\bar\bbeta,
  \qquad
  \Gbar = \Gzero(\thbar),
  \qquad
  \bar{\mathcal{W}} = \mathcal{W}(\Gbar),
  \qquad
  \vbar = \thetahat - \thbar :
\]
\begin{enumerate}[label=\textup{(\roman*)}, leftmargin=*]
  \item \textup{(taxa lenta; nenhuma condição de desenho)}
    \[
      \tfrac12\normn{\Btil\vbar}^{2} + \pen\nGw{\thetahat}
      \;\le\; 3\pen\nGw{\thbar} + 2\normn{\bbar}^{2};
    \]
  \item \textup{(taxa rápida)} se $\gtil > 0$,
    \[
      \normn{\Btil\vbar}^{2} \le \frac{64\pen^{2}\bar{\mathcal{W}}}{\gtil} + 16\normn{\bbar}^{2},
      \quad
      \nGw{\vbar} \le \frac{40\pen\bar{\mathcal{W}}}{\gtil} + \frac{10\normn{\bbar}^{2}}{\pen},
      \quad
      \norm{\vbar}_{2}^{2} \le \frac{64\pen^{2}\bar{\mathcal{W}}}{\gtil^{2}}
      + \frac{16\normn{\bbar}^{2}}{\gtil};
    \]
  \item \textup{(erro de predição)}
    $\normn{\widehat f - f} \le \normn{\Btil\vbar} + 2\normn{\bbar} +
    \normn{\PA\bveps}$, e portanto
    \begin{equation}
      \label{eq:oraculo-min-blocos}
      \normn{\widehat f - f}^{2}
      \;\le\; 3\min_{\bar\bbeta}\Bigl\{\frac{64\pen^{2}\bar{\mathcal{W}}}{\gtil}
      + 20\normn{\bbar}^{2}\Bigr\} + 3\normn{\PA\bveps}^{2},
      \qquad
      \E\bigl[\normn{\PA\bveps}^{2}\mid\bX,\bU\bigr] \le \frac{\sigma^{2}p}{n}.
    \end{equation}
\end{enumerate}
\end{theorem}

Lido no oráculo, $\bar\bbeta = \betastar$, o item (ii) é o Teorema~1(ii) de
E1.5 com $s_{0}$ trocado por $\mathcal{W}(\Gzero(\thetastar))$ e $\norm{\cdot}_{1}$ por
$\nGw{\cdot}$; as constantes são as mesmas. Lido num comparador qualquer, é o
Lema~8 de E1.6 em blocos.

\begin{corollary}[a forma de risco ideal]\label{cor:risco-ideal}
Nas condições do Teorema~\ref{thm:oraculo-blocos}, no evento $\Tw \cap
\{\gtil > 0\}$, seja $\Lambda' \ge \lmaxx(\Gramhat)$ (por exemplo
$\Lambda' = \Lambda$ no evento $\ESig$) e
\[
  \eta \;=\; \frac{1{,}6\,\max_{G}(\pen w_{G})^{2}}{\Lambda'\gtil}.
\]
Então
\begin{equation}
  \label{eq:risco-ideal}
  \normn{\widehat f - f}^{2}
  \;\le\; 120\Lambda'\,\Rid(\thetastar;\eta) + 120\normn{\bb}^{2}
  + 3\normn{\PA\bveps}^{2},
  \quad
  \norm{\thetahat - \thetastar}_{2}^{2}
  \;\le\; \frac{82\Lambda'\,\Rid(\thetastar;\eta) + 64\normn{\bb}^{2}}{\gtil}.
\end{equation}
Com $\pen = \lamG$ (ou $\pen = \pen_{w}$), o Lema~\ref{lem:calib-blocos}(iii)
dá $\eta \le \varrho^{2}\eta_{1}$, com $\eta_{1} = 1{,}6(\lamGone)^{2}/(\Lambda'\gtil)$
o limiar dos pesos~1, e $\Rid(\thetastar;\eta) \le \varrho^{2}\Rid(\thetastar;\eta_{1})$.
\end{corollary}

\begin{remark}[por que não há cone, também em blocos]\label{rem:cone-blocos}
O passo que dispensa o cone é o mesmo de E1.5: a perfilagem do
Lema~\ref{lem:perfil-blocos} tira os níveis do problema, e o autovalor cheio
de E1.4 controla $\nGw{\vbar_{\Gbar}} \le \sqrt{\bar{\mathcal{W}}}\norm{\vbar}_{2} \le
\sqrt{\bar{\mathcal{W}}/\gtil}\,\normn{\Btil\vbar}$ por Cauchy--Schwarz, sobre o espaço
inteiro. Lounici et al.\ (2011) precisam da condição RE em grupos no cone
ponderado (a Hipótese~3.1 deles); aqui ela é dominada pelo autovalor cheio,
como em E1.5, e a constante é a mesma para todo $\Gbar$.
\end{remark}

\begin{remark}[a variante branca, que é o estimador do \texttt{grpreg}]
\label{rem:branca}
O \texttt{grpreg} padroniza as colunas, ortonormaliza cada pedaço e penaliza
com os pesos $\sqrt{|G|}$ as coordenadas ortonormalizadas
(\texttt{08a-sondagem-blocos.md}, §1.3). Na escala original, a penalidade é
$\sum_{G}w_{G}\norm{Q_{G}^{1/2}\btheta_{G}}_{2}$, com $Q_{G} =
\Bmat_{G}\T M_{1}\Bmat_{G}/n$ a Gram centrada do pedaço ($M_{1}$ o centrador).
O Teorema~\ref{thm:oraculo-blocos} vale com qualquer $Q_{G}$ definida positiva
e mensurável em $(\bX,\bU)$, com três trocas: o evento passa a ser
$\{\max_{G}\norm{Q_{G}^{-1/2}(\Btil\T\bveps/n)_{G}}_{2}/w_{G} \le \pen/2\}$ (a
norma dual); $\bar{\mathcal{W}}$ vira $q_{\max}\bar{\mathcal{W}}$ em (ii) e (iii), com
$q_{\max} = \max_{G}\lmaxx(Q_{G})$, porque o Passo~5 da prova lê
$\sum_{G\in\Gbar}w_{G}\norm{Q_{G}^{1/2}\vbar_{G}}_{2} \le
\sqrt{q_{\max}\bar{\mathcal{W}}}\norm{\vbar}_{2}$; e a cota de $\nGw{\vbar}$ fica na
norma com $Q$. O Lema~\ref{lem:calib-blocos}(i) vale com
$Q_{G}^{-1/2}\Psitil_{G}Q_{G}^{-1/2}$ no lugar de $\Psitil_{G}$ (a prova é a
mesma, com $Q_{G}^{-1/2}\Btil_{G}\T/n$ no lugar de $\Btil_{G}\T/n$). Com a
coluna constante $X_{1}\equiv 1$ em $\Amat$, $\PA \succeq P_{1}$, logo
$\MA \preceq M_{1}$ e $\Psitil_{G} \preceq Q_{G}$: o traço é no máximo $|G|$
e a norma, no máximo $1$, e
\[
  \lamz{G} \;\le\; \frac{\sigma}{\sqrt n}\Bigl(\sqrt{|G|} + \sqrt{2\log(|\Gcal|/\alpha)}\Bigr),
\]
\emph{pivotal}, sem $\smax$ nem $\Lambda$; e $q_{\max} \le \lmaxx(\Gramhat)
\le \Lambda$ em $\ESig$. As taxas da §\ref{sec:taxas} passam com a constante
multiplicada por $\Lambda$. É o estimador exato que o \texttt{klopp.balanced}
do piloto roda, com $\pen$ determinístico.
\end{remark}

\section{A aproximação: Besov implica o risco ideal por pedaços}

\begin{lemma}[risco ideal por pedaços sob Besov]\label{lem:risco-pedacos}
Sob as Hipóteses de base e de Besov de E1.3 (em particular $\pi \ge 1$ e
$\seff > 0$, isto é, $\pi(s+1/2) > 1$), seja $b \ge 1$ e $\Gcal$ uma partição
em que, em cada bloco $(\cv,\md)$, os níveis $\lev$ com $2^{\lev} < b$ formam
no máximo um pedaço e cada nível com $2^{\lev} \ge b$ é partido em pedaços
contidos nele com pelo menos $b$ colunas (a forma balanceada é uma delas,
pela Proposição~\ref{prop:balanceada}(i)). Então, para todo $J$ e todo
$\eta > 0$,
\begin{equation}
  \label{eq:risco-pedacos}
  \Rid(\thetastar;\eta) \;\le\; pq\Bigl\{\AG\,C_{g}^{\tau}
  \Bigl(\frac{\eta}{b}\Bigr)^{2s/(2s+1)} b^{(2/\pi-1)_{+}/(2s+1)} + \eta\Bigr\},
  \qquad \tau = \frac{1}{s+1/2},
\end{equation}
com $\AG = 2 + (1 - 2^{1-\pi(s+1/2)})^{-1}$ se $\pi \le 2$ e
$\AG = 2 + (1 - 2^{-2s})^{-1}$ se $\pi \ge 2$ (as duas coincidem em
$\pi = 2$), e $2/\pi = 0$ quando $\pi = \infty$.
\end{lemma}

O expoente é o do Lema~4 do suplemento de Klopp e Pensky (2015, eq.~5.31),
\[
  \textstyle\sum_{l}\min(\norm{a_{l}}^{2},\varepsilon d) \le
  C_{a}^{2/(2r+1)}\varepsilon^{2r/(2r+1)}d^{(2-\nu)_{+}/(\nu(2r+1))},
\]
com
$(d,\varepsilon,r,\nu) = (b,\eta/b,s,\pi)$, mas a hipótese é mais fraca: a
classe~(A3) deles é uma bola $\ell_{\nu}$ ponderada sobre todos os
coeficientes,
$\sum_{k}|a_{k}|^{\nu}(k+1)^{\nu(r+1/2-1/\nu)} \le C_{a}^{\nu}$, que em base de wavelets implica a nossa por nível com $C_{g} = C_{a}$ (no
nível $\lev$, $k+1 \ge 2^{\lev}$), e a nossa, $\Besov{s}{\pi}{\infty}$ em
forma de sequência, não precisa da soma sobre os níveis. É o Lema~9(i) de
E1.6 relido: a mesma contagem por nível, mais um Hölder dentro do pedaço. O
pedaço grosso custa no máximo um $\eta$ por bloco, e é por isso que a forma
balanceada não tem o termo $(x_{+}+1)\eta$ da §3.2 de
\texttt{08a-sondagem-blocos.md}, que vinha de um pedaço por nível grosso.

\section{As taxas}\label{sec:taxas}

\begin{assumption}[regime em blocos]\label{ass:regime-blocos}
Na sequência de problemas da Hipótese~1 de E1.6 ($p$, $q$, $\sigma$, $B_{X}$,
$c_{U}$, $C_{U}$, $\kappa_{1}$, $\kappa_{2}$, $C_{g}$, $s$, $\pi$ e
$\alpha \in (0,1)$ fixos, $J = J_{n} \to \infty$), suponha $2^{J_{n}} \le 2n$,
e tome $\bn = \lceil\log n\rceil$ ($n \ge 3$, de modo que $\bn \ge 2$),
$\Gcal_{n}$ a forma balanceada com $b = \bn$, pesos $w^{(n)}$ com
$\varrho(w^{(n)}) \le \bar\varrho < \infty$ para uma constante $\bar\varrho$ fixa, e
$\pen = \lamG$ do Lema~\ref{lem:calib-blocos}(iii).
\end{assumption}

O WAFC de D44 satisfaz a Hipótese~\ref{ass:regime-blocos} com
$\bar\varrho^{2} = 3$ (Proposição~\ref{prop:balanceada}(iii)); os pesos~1, com
$\bar\varrho = 1$. Duas consequências que as provas usam, provadas na §\ref{sec:provas}:
\begin{equation}
  \label{eq:ordem-lambda}
  (\lamGone)^{2} \;\le\; C_{\pen}\,\frac{\sigma^{2}\bn}{n},
  \qquad
  C_{\pen} = 16\bigl\{2B_{X}^{2}C_{U} + \Lambda\bigl(1 + \tfrac12\log(2pq/\alpha)\bigr)\bigr\},
\end{equation}
e, com $\Acal = \{\smax^{2} \le 2B_{X}^{2}C_{U}\} \cap \ESig$: no evento
$\Acal$, $\pen = \lamG \ge \pen_{w}$ (logo $\Ecal \subseteq \Tw$),
$\gtil > \gamma/2$, $\Amat$ tem posto $p$ e $\lmaxx(\Gramhat) < \Lambda$; e
$\Prob(\Acal) \to 1$ sempre que $pq2^{J_{n}}\log(pq2^{J_{n}})/n \to 0$.

\begin{theorem}[taxa do espaço de aproximação, sem logaritmo]\label{thm:taxa-blocos}
Sob as hipóteses de E1.3, E1.4 e E1.5 e a Hipótese~\ref{ass:regime-blocos},
tome
\begin{equation}
  \label{eq:Jn-blocos}
  J_{n} \;=\; \min\bigl\{J \in \mathbb{N} :\; 2^{J} \ge n^{1/(2\seff+1)}\bigr\}.
\end{equation}
Então:
\begin{enumerate}[label=\textup{(\roman*)}, leftmargin=*]
  \item \textup{(a condição de desenho ao longo de $J_{n}$)}
    $pq2^{J_{n}}\log(pq2^{J_{n}})/n \asymp n^{-2\seff/(2\seff+1)}\log n \to 0$,
    de modo que $\Prob(\Acal) \to 1$;
  \item \textup{(taxa)}
    $\normn{\widehat f - f}^{2} = O_{p}\bigl(n^{-2\seff/(2\seff+1)}\bigr)$;
  \item \textup{(o balanço)} no evento $\Acal \cap \Ecal$, o termo de estimação
    é no máximo $384\bar\varrho^{2}C_{\pen}\sigma^{2}pq(\bn + 2^{J_{n}})/(\gamma n)
    \asymp n^{-2\seff/(2\seff+1)}$, o de viés é $\mathcal{B}_{J_{n}}^{2}
    \asymp 2^{-2J_{n}\seff} \le n^{-2\seff/(2\seff+1)}$, e o de
    $\bn/n$ e o dos níveis, $\sigma^{2}p/n$, são de ordem estritamente menor.
\end{enumerate}
\end{theorem}

\begin{corollary}[componentes e coeficientes funcionais]\label{cor:componentes-blocos}
Nas condições do Teorema~\ref{thm:taxa-blocos}, com
$\widehat g_{\cv\md} = \sum_{\lev<J_{n},\tsl}\coef{\cv}{\md}{\lev}{\tsl}
\wav{\lev}{\tsl}$ e
\[
  \rhoG \;=\; n^{-\seff/(2\seff+1)},
\]
vale $\sum_{\cv,\md}\lVert\widehat g_{\cv\md} -
\gcomp{\cv}{\md}\rVert_{L_{2}[0,1]}^{2} = O_{p}\bigl((\rhoG)^{2}\bigr)$; em
particular $\lVert\widehat g_{\cv\md} - \gcomp{\cv}{\md}\rVert_{L_{2}[0,1]} =
O_{p}(\rhoG)$ para cada $(\cv,\md)$, $\norm{\chat - \bc}_{2} = O_{p}(\rhoG)$,
e $\lVert\widehat{\fcoef{\cv}} - \fcoef{\cv}\rVert_{L_{2}([0,1]^{q})} =
O_{p}(\rhoG)$, $\cv = 1,\dots,p$.
\end{corollary}

\begin{corollary}[taxa adaptativa em blocos]\label{cor:adaptativa-blocos}
Sob as hipóteses de E1.3, E1.4 e E1.5 e a Hipótese~\ref{ass:regime-blocos},
tome $J_{n} = \lceil c\log_{2}n\rceil$ com $c$ na janela
\begin{equation}
  \label{eq:janela-blocos}
  \frac{s}{(2s+1)\seff} \;\le\; c \;<\; 1,
\end{equation}
não vazia se e só se $\seff > s/(2s+1)$, que é a janela do Corolário~5 de
E1.6 (lá escrita $\max\{\tau/2,(2-\tau)/(4\seff)\} \le c < 1$, com
$(2-\tau)/(4\seff) = s/((2s+1)\seff) \ge \tau/2$). Então
\[
  \normn{\widehat f - f}^{2}
  \;=\; O_{p}\Bigl(n^{-2s/(2s+1)}(\log n)^{(2/\pi-1)_{+}/(2s+1)}\Bigr),
\]
e a mesma taxa vale para $\sum_{\cv\md}\lVert\widehat g_{\cv\md} -
\gcomp{\cv}{\md}\rVert_{L_{2}[0,1]}^{2}$. Em particular:
\begin{enumerate}[label=\textup{(\roman*)}, leftmargin=*, nosep]
  \item em $\pi \ge 2$, a taxa é $n^{-2s/(2s+1)}$, \emph{sem logaritmo};
  \item em $\pi < 2$, sobra $(\log n)^{(2/\pi-1)/(2s+1)}$;
  \item o ganho sobre o Corolário~5 de E1.6, $(\log n/n)^{2s/(2s+1)}$, é de
    $(\log n)^{2\seff/(2s+1)}$, na mesma janela de $J_{n}$;
  \item a constante cresce com os pesos no máximo como
    $\bar\varrho^{\,4s/(2s+1)}$ no termo principal.
\end{enumerate}
\end{corollary}

\begin{remark}[onde o logaritmo vai]\label{rem:logaritmo}
O $\pen$ de blocos não é menor que o do LASSO: $(\lamGone)^{2} \asymp
\sigma^{2}\log n/n$, como $\pen_{n}^{2}$ em E1.6. O que muda é o custo por
coordenada: um pedaço fino tem pelo menos $\bn \ge \log|\Gcal_{n}| - O(1)$
coordenadas, e o limiar do risco ideal é $\eta_{n} \asymp \bn/n$, isto é,
$\eta_{n}/\bn \asymp 1/n$ por coordenada, contra $\pen_{n}^{2} \asymp \log n/n$
por coordenada no LASSO. O logaritmo da união sobre as colunas é absorvido
pelo tamanho do pedaço, e é por isso que o pedaço tem de ter pelo menos a
ordem de $\log|\Gcal|$ coordenadas. Em $\pi < 2$ a cota cobra
$\bn^{(2/\pi-1)/(2s+1)}$ das sequências cujos coeficientes grandes estão
isolados, um por pedaço (a ``espiga'' da conferência); um sinal espalhado no
nível não paga nada pelos pedaços.
\end{remark}

\begin{remark}[a cota inferior de Klopp e Pensky em $q = 1$]\label{rem:inferior}
Para $q = 1$, $\bX \perp \bU$ e erro gaussiano, o Teorema~1 de Klopp e Pensky
(2015, eqs.~3.6 e~3.7), com $p$ fixo, dá a cota inferior minimax
$n^{-2r/(2r+1)}$ sobre a classe~(A3) com a base de wavelets periodizadas, que
a (A1)(b) deles admite. Pela observação depois do
Lema~\ref{lem:risco-pedacos}, essa classe com $(r,\nu) = (s,\pi)$ está contida
na hipótese de Besov de E1.3 com $C_{g} = C_{a}$, de modo que a cota inferior
vale sobre a nossa classe. Em $\pi \ge 2$ o
Corolário~\ref{cor:adaptativa-blocos} é, portanto, ótimo a menos de
constante nesse caso particular. Em $\pi < 2$ fica a
$(\log n)^{(2/\pi-1)/(2s+1)}$ dela, o mesmo resto das eqs.~(3.19) e~(3.20) de
Klopp e Pensky com $p$ fixo, e o mesmo expoente da limiarização em blocos de
Cai (1999, Teorema~4, eq.~5.4, p.~908), que vale no modelo de sequência e com
$\alpha \ge 1/p$ (no nosso símbolo, $s \ge 1/\pi$); a nossa condição,
$\seff > s/(2s+1)$, é implicada por ela, porque $s \ge 1/\pi$ dá
$\seff \ge 1/2 > s/(2s+1)$, e é estritamente mais fraca (em $\pi = 1$, é
$s > (1+\sqrt5)/4 \approx 0{,}81$). Que esse resto seja necessário para
este estimador não é provado aqui nem por Klopp e Pensky; para $q \ge 2$ a
cota inferior não existe (E1.9).
\end{remark}

\begin{remark}[o regime pré-assintótico]\label{rem:pre-blocos}
A conferência mede, como E1.6 já tinha medido, que o evento $\ESig$ não vale
nos $n$ do piloto ao longo de $J_{n}$: no desenho da conferência
$\lmaxx(\Gramhat)$ mediano fica entre $2{,}0$ e $2{,}7$ para $n$ de $250$ a
$4000$, contra $\Lambda = 1{,}46$, e a $J = 4$ fixo só entra no evento em
$n = 64000$. O $\lamG$ determinístico continua cobrindo (cobertura $1{,}000$),
porque a calibração é conservadora por um fator perto de $2$, mas o
$\gtil$ realizado é pequeno, e com ele o $\eta$ do
Corolário~\ref{cor:risco-ideal} satura $\Rid$ nesses $n$. As taxas são
assintóticas, como em E1.6.
\end{remark}

\subsection*{Sem condição de desenho: a taxa lenta (E1.13)}

A Proposição~4 de E1.6 é o único enunciado do LASSO que não usa condição de
desenho (D16; no manuscrito, a Proposition~3, ``the unconditional statement of
this paper''), e D47(c) a mantém no corpo em versão de blocos. Ela sai do
Teorema~\ref{thm:oraculo-blocos}(i) com o comparador $\betastar$, que também
não usa $\Gramhat$, e de uma contagem de $\thetastar$ sob a Hipótese de Besov
na norma de grupos com pesos~$1$,
\[
  \nGone{\btheta} \;=\; \sum_{G\in\Gcal}\norm{\btheta_{G}}_{2},
  \qquad
  \vartheta_{\pi} \;=\; \min\Bigl(1 - \frac{1}{\pi},\ \frac12\Bigr)
  \;=\; \frac12 - \Bigl(\frac1\pi - \frac12\Bigr)_{+}
  \qquad (\vartheta_{\infty} = 1/2).
\]
A contagem é a peça nova: o pedaço cobra $\sqrt{\bn}$ em $\lamGone$ e devolve
até $\bn^{\vartheta_{\pi}}$ em $\nGone{\thetastar}$.

Os itens (iii) e (iv) abaixo (E1.14) pedem $2^{J_{n}}$ acima de $n/\log n$,
onde nem $\smax^{2} \le 2B_{X}^{2}C_{U}$ nem $\lmaxx(\Gramhat) \le \Lambda$
valem, e o $\lamG$ determinístico deixa de ser admissível. Com $D = p + d$,
$\xi = \Lambda - \kappa_{2}C_{U} > 0$ e o $c_{\psi} =
7\norm{\psi}_{\infty}^{2}/(12C_{U})$ da Proposição~4(iii) de E1.6, escrevemos
\begin{equation}
  \label{eq:mu-blocos}
  \mu^{\Gcal}_{J} = 1 + c_{\mu}\frac{2^{J}\log(2D)}{n},
  \quad
  c_{\mu} = \max\Bigl\{c_{\psi},\
  \frac{2pB_{X}^{2}(1+qC_{\psi})}{\Lambda}
  \Bigl(\frac{\kappa_{2}C_{U}}{2\xi} + \frac43\Bigr)\Bigr\},
  \quad
  \lamGp = (\mu^{\Gcal}_{J_{n}})^{1/2}\lamG,
\end{equation}
com $C_{\psi}$ o da Proposição~3(iv) de E1.4; quando
$pq2^{J_{n}}\log(pq2^{J_{n}})/n \to 0$, $\lamGp = \lamG\{1 + o(1)\}$.

\setcounter{corollary}{13}
\begin{corollary}[taxa lenta em blocos; nenhuma condição de desenho]\label{cor:lenta-blocos}
Sob as Hipóteses~1 e~2 de E1.5, as hipóteses de E1.3, a
Hipótese~\ref{ass:regime-blocos} e $\Amat$ de posto $p$, suponha
$pq2^{J_{n}}\log(pq2^{J_{n}})/n \to 0$. Então
\begin{equation}
  \label{eq:lenta-blocos}
  \normn{\widehat f - f}^{2}
  \;=\; O_{p}\Bigl(\lamG\nGw{\thetastar} + \mathcal{B}_{J_{n}}^{2}
  + \frac{p}{n}\Bigr),
  \qquad
  \lamG\nGw{\thetastar} \;\le\; \bar\varrho\,\lamGone\nGone{\thetastar}
  \;\le\; \bar\varrho\,\lamGone\norm{\thetastar}_{1},
\end{equation}
com $(\lamGone)^{2} \le C_{\pen}\sigma^{2}\bn/n$ por \eqref{eq:ordem-lambda},
e, com $\jst$ o da forma balanceada em $b = \bn$,
\begin{equation}
  \label{eq:contagem-lenta}
  \nGone{\thetastar} \;\le\; pq\,C_{g}\Bigl\{\frac{1}{1 - 2^{-\seff}}
  + \bn^{-\vartheta_{\pi}}\sum_{\lev=\jst+1}^{J_{n}-1}2^{\lev(1/2-s)}\Bigr\}.
\end{equation}
Em particular:
\begin{enumerate}[label=\textup{(\roman*)}, leftmargin=*]
  \item se $s > 1/2$ e $\seff > 1/4$, então $\nGone{\thetastar} = O(1)$ e,
    para todo $J_{n} \ge \log_{2}n/(4\seff)$ que satisfaça a condição acima
    (existe, porque $\seff > 1/4$),
    $\normn{\widehat f - f}^{2} = O_{p}\bigl((\log n/n)^{1/2}\bigr)$;
  \item se $s < 1/2$ e $\seff > s/2$, então $\nGone{\thetastar} =
    O\bigl(1 + \bn^{-\vartheta_{\pi}}2^{J_{n}(1/2-s)}\bigr)$, e a escolha
    $2^{J_{n}} \asymp (n/\bn^{(2/\pi-1)_{+}})^{1/(1-2s+4\seff)}$ dá
    \[
      \begin{aligned}
      \normn{\widehat f - f}^{2}
      &= O_{p}\Bigl(\Bigl(\frac{\bn^{(2/\pi-1)_{+}}}{n}\Bigr)^{2\seff/(1-2s+4\seff)}\Bigr)\\
      &= O_{p}\Bigl(n^{-2\seff/(1-2s+4\seff)}
        (\log n)^{(2/\pi-1)_{+}\,2\seff/(1-2s+4\seff)}\Bigr);
      \end{aligned}
    \]
    quando $\pi \ge 2$ (isto é, $\seff = s$) a taxa é $n^{-2s/(2s+1)}$,
    \emph{sem logaritmo}: a do Teorema~\ref{thm:taxa-blocos} e do
    Corolário~\ref{cor:adaptativa-blocos}(i), obtida sem nenhuma condição de
    desenho;
  \item \textup{(os pares que (i) e (ii) excluem, todos com $\pi < 2$; E1.14)}
    sem a condição de dimensão acima, com $2^{J_{n}} \le n^{2}$ no lugar de
    $2^{J_{n}} \le 2n$ na Hipótese~\ref{ass:regime-blocos} e $\lamGp$ de
    \eqref{eq:mu-blocos} no lugar de $\lamG$ (no estimador e em
    \eqref{eq:lenta-blocos}), vale \eqref{eq:lenta-blocos} com
    $(\lamGone)^{2} \le (C_{\pen} + 16\Lambda)\sigma^{2}\bn/n$, e:
    \begin{enumerate}[label=\textup{(\alph*)}, leftmargin=*, nosep]
      \item se $s > 1/2$ e $\seff \le 1/4$, a escolha
        $2^{J_{n}} \asymp (n/\log n)^{2/(1+4\seff)}$ dá
        \[
          \normn{\widehat f - f}^{2} =
          O_{p}\bigl((\log n/n)^{4\seff/(1+4\seff)}\bigr),
        \]
        a taxa da Proposição~4(iii)(a) de E1.6;
      \item se $s < 1/2$ e $\seff \le s/2$, a escolha
        $2^{J_{n}} \asymp (n/\bn^{1/\pi})^{1/(1-s+2\seff)}$ dá
        \[
          \normn{\widehat f - f}^{2}
          = O_{p}\Bigl(n^{-2\seff/(1-s+2\seff)}
          (\log n)^{(1/\pi)\,2\seff/(1-s+2\seff)}\Bigr),
        \]
        com $(\log n)^{(1-1/\pi)2\seff/(1-s+2\seff)}$ a menos que a
        Proposição~4(iii)(b) de E1.6;
    \end{enumerate}
  \item \textup{(a fronteira $s = 1/2$, com $\pi > 1$; E1.14)} nas trocas de
    (iii), $\nGone{\thetastar} \le pqC_{g}\{(1-2^{-\seff})^{-1} +
    \bn^{-\vartheta_{\pi}}J_{n}\}$ e: se $\seff \ge 1/4$ ($\pi \ge 4/3$), a
    escolha $2^{J_{n}} \asymp (n/\bn)^{1/(4\seff)}$ dá
    $\normn{\widehat f - f}^{2} = O_{p}\bigl(n^{-1/2}(\log
    n)^{3/2-\vartheta_{\pi}}\bigr)$, isto é, $n^{-1/2}\log n$ em $\pi \ge 2$; se
    $\seff < 1/4$ ($1 < \pi < 4/3$), a escolha $2^{J_{n}} \asymp
    (n/\bn^{1+1/\pi})^{2/(1+4\seff)}$ dá $\normn{\widehat f - f}^{2} =
    O_{p}\bigl(((\log n)^{1+1/\pi}/n)^{4\seff/(1+4\seff)}\bigr)$.
\end{enumerate}
\end{corollary}

\begin{remark}[o que muda em relação à Proposição~4 de E1.6]\label{rem:lenta}
Pela cadeia de \eqref{eq:lenta-blocos} até $\norm{\thetastar}_{1}$, o
Corolário~\ref{cor:lenta-blocos} é a Proposição~4 com a constante
$\bar\varrho$: $(\lamGone)^{2} \asymp \sigma^{2}\log n/n$ é a ordem de
$\pen_{n}^{2} \asymp \sigma^{2}J_{n}/n$, porque lá $J_{n} \asymp \log n$ nos
dois regimes. Nem o pedaço máximo nem o peso máximo mudam o logaritmo: o
pedaço de $2\bn - 1$ colunas entra em $\lamGone$ como $\sqrt{(2\bn-1)/n}$, da
mesma ordem que o termo da união, $\sqrt{\log|\Gcal_{n}|/n}$; e $w_{\max}$ não
aparece, porque $\lamG$ escala como $1/w$ (multiplicar todos os pesos por $c$
divide $\lamG$ por $c$), de modo que só a razão $\varrho$ sobra.

O que essa cadeia joga fora é a contagem \eqref{eq:contagem-lenta}: nos
níveis finos, $\nGone{\thetastar}$ é menor que $\norm{\thetastar}_{1}$ por
$\bn^{-\vartheta_{\pi}}$, e a cota é justa (com a energia do nível espalhada,
em $\pi \ge 2$; com uma espiga por pedaço, em $\pi \le 2$). Em~(i) isso não
muda nada: o pedaço grosso, de norma $O(1)$, domina as duas normas, e a taxa
é $(\log n/n)^{1/2}$ nos dois estimadores. Em~(ii) os níveis finos dominam, e
o ganho $\bn^{-\vartheta_{\pi}}$ cancela o $\sqrt{\bn}$ que o tamanho do pedaço
põe em $\lamGone$, por inteiro em $\pi \ge 2$ e em parte em $\pi < 2$: o
logaritmo da Proposição~4(ii), $(\log n)^{2\seff/(1-2s+4\seff)}$, cai a
$(\log n)^{(2/\pi-1)_{+}2\seff/(1-2s+4\seff)}$. O ganho é
$(\log n)^{\min(1,\,2-2/\pi)\,2\seff/(1-2s+4\seff)}$: o fator
$(\log n)^{2s/(2s+1)}$ inteiro em $\pi \ge 2$, nada em $\pi = 1$. É o mecanismo
da Observação~\ref{rem:logaritmo}: o logaritmo vem do tamanho do pedaço, e o
pedaço o devolve na medida em que o sinal se espalha nele. O paralelo com E1.6
se mantém: lá a Proposição~4(ii) em $\pi \ge 2$ alcança a taxa do Teorema~2;
aqui o Corolário~\ref{cor:lenta-blocos}(ii) alcança a do
Teorema~\ref{thm:taxa-blocos}; nos dois casos sem condição de desenho e só em
$s < 1/2$.

As condições $\seff > 1/4$ em~(i) e $\seff > s/2$ em~(ii) só mordem em
$\pi < 2$, e dizem que o $J_{n}$ do balanço é admissível. São as mesmas de que
a escolha de $J_{n}$ da Proposição~4 de E1.6 precisa para que
$2^{J_{n}}\log d_{n}/n \to 0$, o que aquela prova afirma em todos os regimes:
com $\seff \le 1/4$ em~(i), ou $\seff \le s/2$ em~(ii), o $J_{n}$ de lá tem
$2^{J_{n}} \gtrsim n/\log n$ e o Lema~6 de E1.5 não se aplica (por exemplo,
$(s,\pi) = (0{,}6;\,1)$ e $(0{,}3;\,1{,}5)$).

Os itens (iii) e (iv) (E1.14) cobrem esses pares e a fronteira $s = 1/2$ como
a Proposição~4(iii) e (iv) de E1.6: parar $J_{n}$ no teto daria
$O_{p}((\log n/n)^{2\seff})$, e passar dele com $\lamGp$ dá os expoentes
de (iii), estritamente maiores (Observação~1 de E1.6, que responde por que
nenhuma escolha abaixo do teto faz melhor). O ganho em logaritmo sobre o
LASSO segue o mecanismo de~(ii): nada em (iii)(a), onde o pedaço grosso
domina as duas normas; $(\log n)^{(1-1/\pi)2\seff/(1-s+2\seff)}$ em (iii)(b);
e, em (iv), $(\log n)^{\vartheta_{\pi}}$ se $\seff \ge 1/4$ (meio logaritmo
em $\pi \ge 2$) e $(\log n)^{(1-1/\pi)4\seff/(1+4\seff)}$ se $\seff < 1/4$.
\end{remark}

\begin{remark}[onde o desenho entra, e onde não entra]\label{rem:lenta-desenho}
A cota não assintótica da prova,
\begin{equation}
  \label{eq:lenta-nao-assint}
  \normn{\widehat f - f}^{2} \;\le\; 18\pen\nGw{\thetastar}
  + 24\normn{\bb}^{2} + 3\normn{\PA\bveps}^{2}
  \qquad\text{no evento } \Tw,
\end{equation}
não usa propriedade alguma de $\Gramhat$: nem $\gtil > 0$, nem
$\lmin(\Gramhat)$, nem compatibilidade. Com $\pen \ge \pen_{w}$, calibrado em
$(\bX,\bU)$ pelo Lema~\ref{lem:calib-blocos}(i), ela vale com probabilidade ao
menos $1 - \alpha$ em qualquer desenho com $\Amat$ de posto $p$, inclusive com
$d > n$ ou com a densidade de $\bU$ nula numa parte de $[0,1]^{q}$, onde
$\gtil = 0$ (conferência). O desenho só entra para tornar admissível o $\lamG$
determinístico da Hipótese~\ref{ass:regime-blocos}, e por uma cota
\emph{superior}: o Lema~\ref{lem:calib-blocos}(ii) pede $\smax^{2} \le
2B_{X}^{2}C_{U}$ e $\lmaxx(\Gramhat) \le \Lambda$, e a segunda vem da cauda da
Proposição~3(iv) de E1.4 com $t = \Lambda - \kappa_{2}C_{U}$, que usa só
$\norm{\bZ_{i}}_{2}^{2} \le R_{J}$ e $\lmaxx(\Gram) \le \kappa_{2}C_{U}$ (a
metade superior da Proposição~3(iii)), e não $c_{U}$ nem $\kappa_{1}$;
qualquer $\Lambda > \kappa_{2}C_{U}$ serve. Essa é a diferença para o LASSO:
lá a calibração só pede as normas das colunas ($\smax$); aqui a cauda de Hsu,
Kakade e Zhang tem $\opnorm{\Psitil_{G}}$, e trocá-la pelo traço custaria um
$\sqrt{\log n}$ a mais em $\lamGone$. Em (iii) e (iv), acima do teto, a cota
superior passa a $\lmaxx(\Gramhat) \le \Lambda\mu^{\Gcal}_{J}$, da mesma cauda
com $t$ crescente; continua sem $c_{U}$ nem $\kappa_{1}$.
\end{remark}

\subsection*{Sem condição de desenho, pelo comparador truncado por pedaço (E1.15)}

O Corolário~\ref{cor:lenta-blocos} lê o
Teorema~\ref{thm:oraculo-blocos}(i) só no comparador $\betastar$. Como na
Proposição~8 de E1.6, o comparador truncado por pedaço,
\begin{equation}
  \label{eq:truncado-blocos}
  \thbar_{G} \;=\; \thetastar_{G}\,\mathbf{1}\bigl\{\norm{\thetastar_{G}}_{2} > \pen w_{G}\bigr\},
  \qquad G \in \Gcal,
\end{equation}
cobra na penalidade só os pedaços acima do limiar do pedaço e manda os
outros ao viés. O que sobra é
\begin{equation}
  \label{eq:RG}
  \mathcal{R}_{\Gcal,w}(\btheta;\pen) = \sum_{G\in\Gcal}
  \min\bigl(\pen w_{G}\norm{\btheta_{G}}_{2},\,\norm{\btheta_{G}}_{2}^{2}\bigr),
  \qquad
  \mathcal{R}_{\Gcal,1}(\btheta;t) = \sum_{G\in\Gcal}
  \min\bigl(t\norm{\btheta_{G}}_{2},\,\norm{\btheta_{G}}_{2}^{2}\bigr),
\end{equation}
o par em norma de grupos do risco ideal $\Rid$, com
$\mathcal{R}_{\Gcal,1}(\btheta;t) \le t\nGone{\btheta}$. A contagem nova é a
de $\mathcal{R}_{\Gcal,1}(\thetastar;t)$ sob Besov: nos níveis finos ela junta
a do Corolário~\ref{cor:lenta-blocos} (o ganho $\bn^{-\vartheta_{\pi}}$ da
norma de grupos) com a do Lema~\ref{lem:risco-pedacos} (o corte em que o
pedaço deixa de pagar penalidade e passa a pagar viés).

\begin{corollary}[taxa lenta em blocos pelo comparador truncado; nenhuma condição de desenho]
\label{cor:truncada-blocos}
Sob as Hipóteses~1 e~2 de E1.5, as hipóteses de E1.3, a
Hipótese~\ref{ass:regime-blocos} e $\Amat$ de posto $p$:
\begin{enumerate}[label=\textup{(\roman*)}, leftmargin=*]
  \item \textup{(não assintótico)} no evento $\Tw$, para todo $\pen > 0$, com
    $\thbar$ de \eqref{eq:truncado-blocos},
    \[
      \normn{\widehat f - f}^{2}
      \;\le\; 18\pen\nGw{\thbar} + 48\normn{\Bmat(\thetastar - \thbar)}^{2}
      + 48\normn{\bb}^{2} + 3\normn{\PA\bveps}^{2};
    \]
    se $\pen$ e $w$ são determinísticos,
    $\E\normn{\Bmat(\thetastar - \thbar)}^{2} \le \kappa_{2}C_{U}
    \norm{\thetastar - \thbar}_{2}^{2}$ (e $\le pqB_{X}^{2}C_{U}
    \norm{\thetastar - \thbar}_{2}^{2}$, só com as hipóteses de E1.3),
    \[
      18\pen\nGw{\thbar} + 48\kappa_{2}C_{U}\norm{\thetastar - \thbar}_{2}^{2}
      \;\le\; \max(18,48\kappa_{2}C_{U})\,\mathcal{R}_{\Gcal,w}(\thetastar;\pen),
    \]
    e $\mathcal{R}_{\Gcal,w}(\thetastar;\lamG) \le
    \bar\varrho\,\mathcal{R}_{\Gcal,1}(\thetastar;\lamGone)$. Logo
    \[
      \normn{\widehat f - f}^{2} \;=\; O_{p}\Bigl(\mathcal{R}_{\Gcal,1}(\thetastar;\lamGone)
      + \mathcal{B}_{J_{n}}^{2} + \frac{p}{n}\Bigr)
    \]
    sempre que $pq2^{J_{n}}\log(pq2^{J_{n}})/n \to 0$; com as trocas do
    Corolário~\ref{cor:lenta-blocos}(iii) ($\lamGp$ no estimador, $2^{J_{n}} \le
    n^{2}$), a mesma forma vale com $(\mu^{\Gcal}_{J_{n}})^{1/2}\lamGone$ no lugar
    de $\lamGone$, para toda sequência $J_{n}$.
  \item \textup{(a contagem sob Besov)} Na forma balanceada com $b \ge 2$, para
    todo $J$ e todo $t > 0$, com $\tau = (s+1/2)^{-1}$: se $s < 1/2$,
    \begin{equation}
      \label{eq:contagem-truncada-blocos}
      \begin{aligned}
      \mathcal{R}_{\Gcal,1}(\thetastar;t) \le pq\Bigl\{&\frac{C_{g}\,t}{1 - 2^{-\seff}}\\
      &+ C_{g}^{\tau}\,t^{4s/(2s+1)}\,b^{((2/\pi-1)_{+}-2s)/(2s+1)}
      \Bigl(\frac{1}{1 - 2^{s-1/2}} + \frac{1}{1 - 2^{-\min(\pi,2)\seff}}\Bigr)\Bigr\};
      \end{aligned}
    \end{equation}
    se $s \ge 1/2$, nada se ganha em ordem: $\mathcal{R}_{\Gcal,1}(\thetastar;t)
    \le t\nGone{\thetastar}$, e para $t \le C_{g}$ a classe contém um
    $\thetastar$ com $\mathcal{R}_{\Gcal,1}(\thetastar;t) \ge tC_{g}$ (se
    $s > 1/2$) e um com $\mathcal{R}_{\Gcal,1}(\thetastar;t) \ge
    tC_{g}b^{-\vartheta_{\pi}}/2$ por nível fino $\lev < J$ com
    $2^{\lev} \le C_{g}b^{1-\vartheta_{\pi}}/t$ (se $s = 1/2$), da ordem de
    \eqref{eq:contagem-lenta}.
  \item \textup{($s < 1/2$ e $\seff > s/(2s+1)$)} com $J_{n} = \lceil
    c\log_{2}n\rceil$ e $c$ na janela \eqref{eq:janela-blocos},
    \[
      \normn{\widehat f - f}^{2}
      \;=\; O_{p}\Bigl(n^{-2s/(2s+1)}(\log n)^{(2/\pi-1)_{+}/(2s+1)}\Bigr),
    \]
    a taxa do Corolário~\ref{cor:adaptativa-blocos}, sem condição de desenho;
  \item \textup{($s < 1/2$ e $\seff \le s/(2s+1)$, só com $\pi < 2$)} com as
    trocas de (i) e a escolha $2^{J_{n}} \asymp (n^{2s}\bn^{-(2s+2/\pi-1)/2})^{1/(s+\seff(2s+1))}$,
    que passa de $n/\bn$,
    \[
      \normn{\widehat f - f}^{2}
      \;=\; O_{p}\Bigl(n^{-e}(\log n)^{e(2s+2/\pi-1)/(4s)}\Bigr),
      \qquad e = \frac{4s\seff}{s+\seff(2s+1)};
    \]
    em $\seff = s/(2s+1)$, $n^{-2s/(2s+1)}(\log n)^{(2s+2/\pi-1)/(2(2s+1))}$.
\end{enumerate}
\end{corollary}

\begin{remark}[o que muda diante do Corolário~\ref{cor:lenta-blocos}]\label{rem:truncada-blocos}
\emph{Em $s < 1/2$ e $\pi < 2$, a taxa melhora em todo par, com os mesmos
expoentes de $n$ da Proposição~8 de E1.6.} O item (iii) dá, na janela do
Corolário~\ref{cor:adaptativa-blocos}, a taxa dele, sem a condição de desenho
de E1.4: em $s < 1/2$, a taxa de predição do Corolário~11 não precisa da
condição de desenho (a das componentes continua precisando). O item (iv) troca o Corolário~\ref{cor:lenta-blocos}(iii)(b) e a
parte de (ii) com $s/2 < \seff \le s/(2s+1)$. Em $\pi \ge 2$ a taxa,
$n^{-2s/(2s+1)}$ sem logaritmo, é a do Corolário~\ref{cor:lenta-blocos}(ii),
mas qualquer $c \in [1/(2s+1),1)$ serve.

\emph{Em $s \ge 1/2$ nada muda}, pela razão do LASSO: o pedaço grosso, de
norma $C_{g}$, custa $\pen C_{g}$ na penalidade ou $C_{g}^{2}$ no viés, e a
rota lenta não desce de $\lamGone \asymp (\log n/n)^{1/2}$. Os itens (i),
(iii)(a) e (iv) do Corolário~\ref{cor:lenta-blocos} continuam sendo o
enunciado sem condição de desenho nesses pares.

\emph{Diante do LASSO truncado}, o ganho no logaritmo é o mesmo mecanismo da
Observação~\ref{rem:logaritmo}: $(\log n)^{2\seff/(2s+1)}$ em (iii), o ganho
do Corolário~\ref{cor:adaptativa-blocos} sobre o Corolário~5 de E1.6, e
$(\log n)^{e\seff/(2s)}$ em (iv).

\emph{O desenho.} Como na Observação~\ref{rem:lenta-desenho}, só entram cotas
superiores: a de $\lmaxx(\Gramhat)$ na calibração ($\Lambda$, ou
$\Lambda\mu^{\Gcal}_{J}$ acima do teto) e, no viés novo,
$\lmaxx(\Gram) \le \kappa_{2}C_{U}$ (ou $pqB_{X}^{2}C_{U}$); nem $c_{U}$, nem
$\kappa_{1}$, nem $\gtil$. Os pesos custam o fator $\bar\varrho$, como no
Corolário~\ref{cor:lenta-blocos}. O enunciado é de predição: as componentes
continuam pedindo $\lmin(\Gramhat)$ (Corolário~\ref{cor:componentes-blocos}).
\end{remark}

\section{Diante do Teorema~2 de Klopp e Pensky (2015)}\label{sec:kp}

Lido no \emph{Annals} (L7): a norma (3.1) e o estimador (3.2)--(3.3) na
p.~1280, as condições (3.13) e a escolha (3.14) de $\widehat\delta$ na
p.~1283, o Teorema~2 e a Observação~2 nas pp.~1283--1284, a (A6) na p.~1278.

\begin{itemize}[leftmargin=*]
  \item \textbf{O estimador.} Os pedaços de Klopp e Pensky têm
    $d \approx \log n$ coeficientes consecutivos no índice da base,
    atravessando níveis, com pesos~$1$, e o coeficiente da constante é um
    pedaço unitário penalizado. Aqui os níveis $\lvl{\cv}$ são livres (D3), os
    pedaços finos ficam dentro de um nível e os níveis grossos formam um
    pedaço por bloco, e os pesos são de razão limitada, o que cobre os
    $\sqrt{|G|}$ do \texttt{grpreg} (Proposição~\ref{prop:balanceada}(iii)) e
    a variante branca (Observação~\ref{rem:branca}).
  \item \textbf{O $\pen$.} O Teorema~2 deles usa $\delta = 2\widehat\delta$,
    isto é, $\pen = \widehat\delta$, com $\widehat\delta \asymp
    \sqrt{\log p/n}$ e a condição (3.16), $p^{\mu} \ge 2n$, que com $p$ fixo
    força $\mu\log p \ge \log(2n)$: o $\log p$ do enunciado esconde um
    $\log n$, e $\widehat\delta \asymp \sigma\sqrt{\log n/n}$, a ordem de
    $\lamGone$ em \eqref{eq:ordem-lambda}.
  \item \textbf{$p$ fixo.} A (A6), $(s+s_{0})(1+\log n) \le p$, usada na
    prova do Teorema~2 (Lema~3 do suplemento deles), é de alta dimensão e
    falha com $p$ fixo e $n$ grande, que é o regime da
    Hipótese~\ref{ass:regime-blocos}. Nada aqui a usa.
  \item \textbf{O viés e a dimensão.} A (3.13) exige $L + 1 = n^{\varsigma}$,
    $1/2 \le \varsigma < 1$, e $r^{*} = \min_{j}(r_{j}\wedge r'_{j}) \ge
    r_{0}^{*} > (2\varsigma)^{-1}$; pela Observação~2, $\varsigma = 1/2$ pede
    $r^{*} \ge 2$, e $\varsigma = 3/4$ pede $r^{*} \ge 1$ ao preço de
    $n \ge N_{3}^{2}$. A razão é que o viés de truncamento,
    $C_{a}^{2}s\,n^{-2r^{*}\varsigma}$, é levado à ordem paramétrica
    $C_{a}^{2}s/n$, e o termo de viés no ajuste é controlado pela norma dual,
    como o de ruído (Lema~2 do suplemento deles, eq.~5.25). Aqui o viés entra
    por Young (Passo~2 da prova do Teorema~\ref{thm:oraculo-blocos}), sem
    norma dual, e só precisa não dominar a taxa: a janela
    \eqref{eq:janela-blocos} pede $\seff > s/(2s+1)$, que é sempre menor que
    $1/2$, enquanto $r^{*} > (2\varsigma)^{-1} > 1/2$ (com $r^{*} = \seff$ em
    $\pi < 2$ e $r^{*} = s$ em $\pi \ge 2$). Os pares com
    $s/(2s+1) < \seff \le 1/2$ estão no Corolário~\ref{cor:adaptativa-blocos}
    e fora do Teorema~2 deles; entre eles está o cenário não homogêneo de D27,
    com $\seff = 1/2$ ($\pi = 1$, $s = 1$). A nossa dimensão também é
    polinomial, $2^{J_{n}} \asymp n^{c}$, mas com $c \ge s/((2s+1)\seff)$,
    que fica abaixo de $1/2$ quando $\pi \ge 2$ e $s > 1/2$.
  \item \textbf{A taxa.} Com $p$ fixo e $\log p \asymp \log n$, a (3.19)
    deles é $n^{-2r/(2r+1)}(\log n)^{(2-\nu)_{+}/(\nu(2r+1))}$, exatamente a do
    Corolário~\ref{cor:adaptativa-blocos} com $(r,\nu) = (s,\pi)$. O
    Corolário~1 deles, eq.~(3.21), imprime $2(2-\nu)_{+}/(\nu(2r+1))$, com um
    fator $2$ que (3.6), (3.19) e (3.20) não sustentam
    (\texttt{08a-sondagem-blocos.md}, §4.3, conferido em 2026-10-02).
  \item \textbf{O que eles têm e este arquivo não.} A cota inferior
    (Teorema~1, eqs.~3.6 e~3.7), que aqui só vale emprestada para $q = 1$ e
    $\bX\perp\bU$ (Observação~\ref{rem:inferior}); o regime de alta dimensão,
    $p \gg n$, com a isometria restrita em média (A5) e a Gram restrita do
    Lema~1 deles, onde aqui se usa o autovalor cheio de E1.4, que exige
    $pq2^{J}\log(pq2^{J})/n \to 0$; e a versão em risco quadrático médio
    (Teorema~3). O que este arquivo tem e eles não: $q \ge 2$ moduladoras com o
    termo cruzado, $\bX$ dependente de $\bU$ (D13), os níveis livres, pesos de
    razão limitada, e as três condições acima ($p$ fixo sem (A6), viés sem
    norma dual, sem $r^{*} > (2\varsigma)^{-1}$).
\end{itemize}

\section{Provas}\label{sec:provas}

\subsection*{Proposição~\ref{prop:balanceada}}

Todos os blocos têm a mesma estrutura, de modo que basta olhar um. O pedaço
grosso tem $\sum_{\lev=0}^{\min(\jst,J-1)}2^{\lev} = 2^{\min(\jst,J-1)+1}-1$
colunas. Se $J \le \jst+1$, são $2^{J}-1$ e não há nível fino. Caso
contrário, $2^{\jst} < b \le 2^{\jst+1}$ e, $b$ sendo inteiro,
$2^{\jst} \le b-1$; logo $b - 1 \le 2^{\jst+1}-1 \le 2b-3$. Num nível fino há
$k_{\lev} = \lfloor 2^{\lev}/b\rfloor \ge 1$ pedaços: $k_{\lev}-1$ com $b$
colunas e o último com $b + r_{\lev}$, $r_{\lev} = 2^{\lev} - k_{\lev}b \in
\{0,\dots,b-1\}$; todos entre $b$ e $2b-1$, o que dá (i). Para (ii), um bloco
tem $1 + \sum_{\lev=\jst+1}^{J-1}\lfloor 2^{\lev}/b\rfloor \le 1 +
(2^{J}-2^{\jst+1})/b \le 1 + 2^{J}/b$ pedaços. Para (iii), com
$w_{G} = \sqrt{|G|}$, $\varrho^{2} = \max|G|/\min|G|$; sem nível fino, todos os
pedaços têm $2^{J}-1$ colunas e $\varrho = 1$; com nível fino, por (i),
$\min|G| \ge b-1$ e $\max|G| \le 2b-1$, e $(2b-1)/(b-1) = 2 + 1/(b-1) \le 3$.
Os valores em $b = 6$ e $7$ são os de \texttt{08a-sondagem-blocos.md},
§11.2: o pedaço grosso tem $7$ colunas e os finos, de $6$ a $10$ ($b = 6$) e
de $7$ a $11$ ($b = 7$). \qed

\subsection*{Lema~\ref{lem:calib-blocos}}

\emph{(i).} Condicione em $(\bX,\bU)$. Nessa condicional $\Btil$ é constante e,
pelas Hipóteses~1 e~2 de E1.5, os $\varepsilon_{i}$ são independentes com
$\E[e^{a_{i}\varepsilon_{i}}\mid\bX,\bU] \le e^{\sigma^{2}a_{i}^{2}/2}$, logo
$\E[e^{a\T\bveps}\mid\bX,\bU] \le e^{\sigma^{2}\norm{a}_{2}^{2}/2}$ para todo
$a \in \R^{n}$: é a hipótese (2.1) do Teorema~2.1 de Hsu, Kakade e Zhang
(2012), com $\mu = 0$. Fixe $G$ e escreva $M_{G} = \Btil_{G}\T/n \in
\R^{|G|\times n}$, de modo que $(\Btil\T\bveps/n)_{G} = M_{G}\bveps$, e
$\Sigma_{G} = M_{G}\T M_{G} \in \R^{n\times n}$. O teorema publicado toma uma
matriz quadrada $A \in \R^{n\times n}$ com $\Sigma = A\T A$; tome
$A = \Sigma_{G}^{1/2}$, que é quadrada e tem o mesmo $\Sigma_{G}$, e note que
$\norm{\Sigma_{G}^{1/2}\bveps}_{2}^{2} = \bveps\T\Sigma_{G}\bveps =
\norm{M_{G}\bveps}_{2}^{2}$: a forma retangular só entra por
$\Sigma_{G} = M_{G}\T M_{G}$. Como $\operatorname{tr}\Sigma_{G} =
\operatorname{tr}(M_{G}M_{G}\T) = \operatorname{tr}\Psitil_{G}/n$,
$\opnorm{\Sigma_{G}} = \opnorm{M_{G}M_{G}\T} = \opnorm{\Psitil_{G}}/n$ e
$\operatorname{tr}(\Sigma_{G}^{2}) \le
\opnorm{\Sigma_{G}}\operatorname{tr}\Sigma_{G}$, o teorema dá, para todo
$t > 0$, com probabilidade condicional ao menos $1 - e^{-t}$,
\[
  \begin{aligned}
  \norm{(\Btil\T\bveps/n)_{G}}_{2}^{2}
  &\le \sigma^{2}\Bigl(\operatorname{tr}\Sigma_{G} +
  2\sqrt{\operatorname{tr}(\Sigma_{G}^{2})\,t} + 2\opnorm{\Sigma_{G}}t\Bigr)
  \le \sigma^{2}\Bigl(\sqrt{\operatorname{tr}\Sigma_{G}} +
  \sqrt{2\opnorm{\Sigma_{G}}t}\Bigr)^{2}\\
  &= \frac{\sigma^{2}}{n}\Bigl(\sqrt{\operatorname{tr}\Psitil_{G}} +
  \sqrt{2\opnorm{\Psitil_{G}}t}\Bigr)^{2},
  \end{aligned}
\]
porque $2\sqrt{\opnorm{\Sigma_{G}}\operatorname{tr}\Sigma_{G}\,t} \le
2\sqrt{2\opnorm{\Sigma_{G}}\operatorname{tr}\Sigma_{G}\,t}$. Com
$t = \log(|\Gcal|/\alpha)$, cada pedaço falha com probabilidade condicional no
máximo $\alpha/|\Gcal|$; a união sobre os $|\Gcal|$ pedaços e a torre dão
$\Prob(\Ecal) \ge 1-\alpha$. Em $\Ecal$, para todo $G$,
$\norm{(\Btil\T\bveps/n)_{G}}_{2}/w_{G} \le \lamz{G}/w_{G} \le \pen_{w}/2 \le
\pen/2$, logo $\Ecal \subseteq \Tw$.

\emph{(ii).} $\MA \preceq I_{n}$ dá $\Psitil_{G} = \Bmat_{G}\T\MA\Bmat_{G}/n
\preceq \Bmat_{G}\T\Bmat_{G}/n = \Gramhat_{GG}$, a submatriz principal de
$\Gramhat$ em $G$. Logo $\operatorname{tr}\Psitil_{G} \le
\sum_{a\in G}(\Gramhat)_{aa} \le |G|\smax^{2}$ e
$\opnorm{\Psitil_{G}} \le \lmaxx(\Gramhat_{GG}) \le \lmaxx(\Gramhat)$
(entrelaçamento de Cauchy). No evento do enunciado,
$\smax \le B_{X}\sqrt{2C_{U}}$ e $\lmaxx(\Gramhat) \le \Lambda$.

\emph{(iii).} Para todo $G$,
\[
  \lamG w_{G} = 2\max_{G'}\bigl(\lamzbar{G'}w_{G}/w_{G'}\bigr) \le
  2\max_{G'}\lamzbar{G'}\cdot\max_{G'}(w_{G}/w_{G'}) \le \varrho\lamGone,
\]
e
$(\lamG)^{2}\mathcal{W}(\Hcal) = \sum_{G\in\Hcal}(\lamG w_{G})^{2} \le
\varrho^{2}(\lamGone)^{2}|\Hcal|$. A mesma conta vale com os $\lamz{G}$. \qed

\subsection*{Teorema~\ref{thm:oraculo-blocos}}

Fixe um comparador e escreva $\norm{\cdot}$ por $\nGw{\cdot}$ nesta prova.

\emph{Passo 1 (desigualdade básica).} Pelo Lema~\ref{lem:perfil-blocos}(i),
$\thetahat$ minimiza o objetivo residualizado; comparando com $\thbar$,
\[
  \normn{\widetilde\bY - \Btil\thetahat}^{2} + 2\pen\norm{\thetahat}
  \;\le\; \normn{\widetilde\bY - \Btil\thbar}^{2} + 2\pen\norm{\thbar}.
\]
Como $\bY = \Amat\bar\bc + \Bmat\thbar + \bbar + \bveps$ e $\MA\Amat = 0$, vale
$\widetilde\bY = \Btil\thbar + r$ com $r = \MA(\bbar + \bveps)$, donde
\begin{equation}
  \label{eq:basica-blocos}
  \normn{\Btil\vbar}^{2} + 2\pen\norm{\thetahat}
  \;\le\; 2\inner{r}{\Btil\vbar}_{n} + 2\pen\norm{\thbar}.
\end{equation}

\emph{Passo 2 (os dois ruídos).} Como $\MA$ é simétrica e $\MA\Btil = \Btil$,
$\inner{r}{\Btil\vbar}_{n} = \inner{\bbar + \bveps}{\Btil\vbar}_{n}$. O termo
estocástico, pela norma dual de $\norm{\cdot}$, que é
$\max_{G}\norm{z_{G}}_{2}/w_{G}$: em $\Tw$,
\[
  2\bigl|\inner{\bveps}{\Btil\vbar}_{n}\bigr|
  = 2\Bigl|\sum_{G}\inner{(\Btil\T\bveps/n)_{G}}{\vbar_{G}}\Bigr|
  \le 2\max_{G}\frac{\norm{(\Btil\T\bveps/n)_{G}}_{2}}{w_{G}}
  \sum_{G}w_{G}\norm{\vbar_{G}}_{2}
  \le \pen\norm{\vbar}.
\]
O termo determinístico, por Cauchy--Schwarz e Young, como em E1.5:
\[
  2|\inner{\bbar}{\Btil\vbar}_{n}| = 2|\inner{\MA\bbar}{\Btil\vbar}_{n}| \le
  2\normn{\bbar}\normn{\Btil\vbar} \le \tfrac12\normn{\Btil\vbar}^{2} +
  2\normn{\bbar}^{2}.
\]
Levando os dois a \eqref{eq:basica-blocos},
\begin{equation}
  \label{eq:basica2-blocos}
  \tfrac12\normn{\Btil\vbar}^{2} + 2\pen\norm{\thetahat}
  \;\le\; \pen\norm{\vbar} + 2\pen\norm{\thbar} + 2\normn{\bbar}^{2}.
\end{equation}

\emph{Passo 3 (taxa lenta).} $\norm{\vbar} \le \norm{\thetahat} +
\norm{\thbar}$ em \eqref{eq:basica2-blocos} dá (i). Nenhuma propriedade de
$\Gramhat$ foi usada.

\emph{Passo 4 (decomponibilidade).} Para $G \notin \Gbar$, $\thbar_{G} = 0$ e
$\thetahat_{G} = \vbar_{G}$; para $G \in \Gbar$,
$\norm{\thetahat_{G}}_{2} \ge \norm{\thbar_{G}}_{2} - \norm{\vbar_{G}}_{2}$.
Logo $\norm{\thetahat} \ge \norm{\thbar} - \norm{\vbar_{\Gbar}} +
\norm{\vbar_{\Gbar^{c}}}$, e $\norm{\vbar} = \norm{\vbar_{\Gbar}} +
\norm{\vbar_{\Gbar^{c}}}$. Substituindo em \eqref{eq:basica2-blocos} e
cancelando $2\pen\norm{\thbar}$,
\begin{equation}
  \label{eq:cone-blocos}
  \tfrac12\normn{\Btil\vbar}^{2} + \pen\norm{\vbar_{\Gbar^{c}}}
  \;\le\; 3\pen\norm{\vbar_{\Gbar}} + 2\normn{\bbar}^{2}.
\end{equation}

\emph{Passo 5 (os dois casos, sem cone).} Por Cauchy--Schwarz e pela
definição de $\gtil$,
\begin{equation}
  \label{eq:cs-blocos}
  \norm{\vbar_{\Gbar}} = \sum_{G\in\Gbar}w_{G}\norm{\vbar_{G}}_{2}
  \le \Bigl(\sum_{G\in\Gbar}w_{G}^{2}\Bigr)^{1/2}
  \Bigl(\sum_{G\in\Gbar}\norm{\vbar_{G}}_{2}^{2}\Bigr)^{1/2}
  \le \sqrt{\bar{\mathcal{W}}}\,\norm{\vbar}_{2}
  \le \sqrt{\bar{\mathcal{W}}/\gtil}\;\normn{\Btil\vbar},
\end{equation}
porque $\normn{\Btil\vbar}^{2} = \vbar\T\GBA\vbar \ge \gtil\norm{\vbar}_{2}^{2}$.
\textbf{Caso 1:} $2\normn{\bbar}^{2} \le \pen\norm{\vbar_{\Gbar}}$. Então
\eqref{eq:cone-blocos} dá $\frac12\normn{\Btil\vbar}^{2} +
\pen\norm{\vbar_{\Gbar^{c}}} \le 4\pen\norm{\vbar_{\Gbar}}$, e com
\eqref{eq:cs-blocos}, $\frac12\normn{\Btil\vbar}^{2} \le
4\pen\sqrt{\bar{\mathcal{W}}/\gtil}\,\normn{\Btil\vbar}$, isto é
$\normn{\Btil\vbar} \le 8\pen\sqrt{\bar{\mathcal{W}}/\gtil}$ e
$\normn{\Btil\vbar}^{2} \le 64\pen^{2}\bar{\mathcal{W}}/\gtil$. Do mesmo display,
$\norm{\vbar_{\Gbar^{c}}} \le 4\norm{\vbar_{\Gbar}}$, logo $\norm{\vbar} \le
5\norm{\vbar_{\Gbar}} \le 5\sqrt{\bar{\mathcal{W}}/\gtil}\cdot 8\pen\sqrt{\bar{\mathcal{W}}/\gtil}
= 40\pen\bar{\mathcal{W}}/\gtil$; e $\norm{\vbar}_{2}^{2} \le
\normn{\Btil\vbar}^{2}/\gtil \le 64\pen^{2}\bar{\mathcal{W}}/\gtil^{2}$.
\textbf{Caso 2:} $2\normn{\bbar}^{2} > \pen\norm{\vbar_{\Gbar}}$. Então
\eqref{eq:cone-blocos} dá $\frac12\normn{\Btil\vbar}^{2} <
8\normn{\bbar}^{2}$, isto é $\normn{\Btil\vbar}^{2} < 16\normn{\bbar}^{2}$;
também $\pen\norm{\vbar_{\Gbar^{c}}} < 8\normn{\bbar}^{2}$ e
$\pen\norm{\vbar_{\Gbar}} < 2\normn{\bbar}^{2}$, donde $\norm{\vbar} <
10\normn{\bbar}^{2}/\pen$; e $\norm{\vbar}_{2}^{2} \le
\normn{\Btil\vbar}^{2}/\gtil < 16\normn{\bbar}^{2}/\gtil$. Cada quantidade é
majorada pelo máximo dos dois casos e o máximo pela soma: é (ii).

\emph{Passo 6 (erro de predição).} Por \eqref{eq:decomp-blocos},
$\widehat f - f = \Btil\vbar + \PA(\bbar + \bveps) - \bbar$, e
$\normn{\PA\bbar} \le \normn{\bbar}$; a desigualdade triangular dá a primeira
parte de (iii), e $(x+y+z)^{2} \le 3(x^{2}+y^{2}+z^{2})$ com (ii) dá
$\normn{\widehat f - f}^{2} \le 3(64\pen^{2}\bar{\mathcal{W}}/\gtil + 16\normn{\bbar}^{2}
+ 4\normn{\bbar}^{2} + \normn{\PA\bveps}^{2})$. Como $\Tw$, $\gtil$ e $\pen$ não
dependem do comparador, as cotas valem simultaneamente para todos, e
pode-se tomar o mínimo. A cota de $\E[\normn{\PA\bveps}^{2}\mid\bX,\bU]$ é a
do Passo~6 de E1.5: $\operatorname{tr}(\PA) \le p$. \qed

\subsection*{Corolário~\ref{cor:risco-ideal}}

Tome $T = \{G : \norm{\thetastar_{G}}_{2}^{2} > \eta\}$ e o comparador
$\bar\bbeta = (\bc,\thbar)$ com $\thbar_{G} = \thetastar_{G}\mathbf{1}\{G\in
T\}$. Então $\Gbar \subseteq T$ e, pela definição de $\eta$,
$\pen^{2}\bar{\mathcal{W}} \le \sum_{G\in T}(\pen w_{G})^{2} \le
|T|\max_{G}(\pen w_{G})^{2} = |T|\,\eta\Lambda'\gtil/1{,}6$. O resíduo é
$\bbar = \bb + \Bmat(\thetastar - \thbar)$, e como $\Bmat\T\Bmat/n$ é submatriz
principal de $\Gramhat$,
\[
  \normn{\bbar}^{2} \le 2\normn{\bb}^{2} + 2\lmaxx(\Gramhat)
  \norm{\thetastar - \thbar}_{2}^{2}
  \le 2\normn{\bb}^{2} + 2\Lambda'\sum_{G\notin T}\norm{\thetastar_{G}}_{2}^{2}.
\]
Por \eqref{eq:oraculo-min-blocos},
$\normn{\widehat f - f}^{2} \le 192\pen^{2}\bar{\mathcal{W}}/\gtil + 60\normn{\bbar}^{2} +
3\normn{\PA\bveps}^{2} \le 120\Lambda'\{\eta|T| + \sum_{G\notin
T}\norm{\thetastar_{G}}^{2}\} + 120\normn{\bb}^{2} + 3\normn{\PA\bveps}^{2}$,
e o termo entre chaves é exatamente $\Rid(\thetastar;\eta)$. Para as
componentes, o Teorema~\ref{thm:oraculo-blocos}(ii) dá
$\norm{\vbar}_{2}^{2} \le 64\pen^{2}\bar{\mathcal{W}}/\gtil^{2} +
16\normn{\bbar}^{2}/\gtil \le (40\Lambda'\eta|T| + 32\normn{\bb}^{2} +
32\Lambda'\sum_{G\notin T}\norm{\thetastar_{G}}^{2})/\gtil \le
(40\Lambda'\Rid + 32\normn{\bb}^{2})/\gtil$, e
$\norm{\thetahat - \thetastar}_{2}^{2} \le 2\norm{\vbar}_{2}^{2} +
2\norm{\thbar - \thetastar}_{2}^{2} \le (80\Lambda'\Rid +
64\normn{\bb}^{2})/\gtil + 2\Rid$; como $\gtil \le \lmin(\Gramhat_{\Bmat\Bmat})
\le \lmaxx(\Gramhat) \le \Lambda'$ (o complemento de Schur é dominado pelo
bloco), $2\Rid \le 2\Lambda'\Rid/\gtil$, e sai \eqref{eq:risco-ideal}. A
última afirmação é o Lema~\ref{lem:calib-blocos}(iii) e as duas propriedades
de $\Rid$ da §1. \qed

\subsection*{Lema~\ref{lem:risco-pedacos}}

Fixe um bloco e escreva $\theta_{\lev} = \theta^{*}_{\cv\md,\lev\cdot} \in
\R^{2^{\lev}}$ (nulo para $\lev \ge J$, o que só diminui o risco),
$\Gcal_{\lev}$ para os pedaços do nível fino $\lev$ e
$r_{\lev} = \sum_{G\in\Gcal_{\lev}}\min(\norm{\theta_{G}}_{2}^{2},\eta)$.
O pedaço grosso, se houver, contribui no máximo $\eta$; logo a contribuição do
bloco é no máximo $\eta + \sum_{\lev\text{ fino}}r_{\lev}$. A primeira cota de
$r_{\lev}$ é a contagem: cada pedaço de $\Gcal_{\lev}$ tem pelo menos $b$ das
$2^{\lev}$ colunas do nível, logo $r_{\lev} \le |\Gcal_{\lev}|\eta \le
2^{\lev}\eta/b$.

\emph{Caso $\pi \le 2$.} Escreva $a = \pi(s+1/2) > 1$. Para $u \ge 0$,
$\min(u,\eta) \le \eta^{1-\pi/2}u^{\pi/2}$ (se $u \le \eta$,
$u = u^{\pi/2}u^{1-\pi/2} \le u^{\pi/2}\eta^{1-\pi/2}$; se $u > \eta$,
$\eta = \eta^{\pi/2}\eta^{1-\pi/2} < u^{\pi/2}\eta^{1-\pi/2}$), e
$\norm{x}_{2} \le \norm{x}_{\pi}$ para $\pi \le 2$. Logo a segunda cota é
\[
  r_{\lev} \le \eta^{1-\pi/2}\sum_{G\in\Gcal_{\lev}}\norm{\theta_{G}}_{\pi}^{\pi}
  \le \eta^{1-\pi/2}\norm{\theta_{\lev}}_{\pi}^{\pi}
  \le \eta^{1-\pi/2}C_{g}^{\pi}2^{-\lev(a-1)},
\]
pela Hipótese de Besov, porque os pedaços de $\Gcal_{\lev}$ estão contidos no
nível $\lev$ e são disjuntos. Defina $x \in \R$ por
$2^{xa} = bC_{g}^{\pi}\eta^{-\pi/2}$, de modo que as duas cotas se igualam em
$\lev = x$, com valor $h = 2^{x}\eta/b = \eta^{1-\pi/2}C_{g}^{\pi}2^{-x(a-1)}$.
Somando a primeira sobre os níveis $\lev \le x$,
$\sum_{\lev\le x}2^{\lev}\eta/b \le 2^{\lfloor x\rfloor+1}\eta/b \le 2h$ (soma
vazia se $x < 0$); somando a segunda sobre os $\lev > x$, como o primeiro
inteiro $\lev > x$ tem $2^{-\lev(a-1)} \le 2^{-x(a-1)}$,
$\sum_{\lev>x}\eta^{1-\pi/2}C_{g}^{\pi}2^{-\lev(a-1)} \le
h\sum_{i\ge0}2^{-i(a-1)} = h/(1-2^{1-a})$. A contribuição do bloco é, então,
no máximo $\eta + \{2 + (1-2^{1-a})^{-1}\}h$. Resta calcular $h$: como
$1/a = \tau/\pi$, $2^{x} = b^{\tau/\pi}C_{g}^{\tau}\eta^{-\tau/2}$ e
$h = b^{\tau/\pi-1}C_{g}^{\tau}\eta^{1-\tau/2}$; com $1 - \tau/2 =
2s/(2s+1)$,
\[
  h = C_{g}^{\tau}\Bigl(\frac{\eta}{b}\Bigr)^{2s/(2s+1)}
  b^{\tau/\pi - 1 + 2s/(2s+1)}
  = C_{g}^{\tau}\Bigl(\frac{\eta}{b}\Bigr)^{2s/(2s+1)}
  b^{\tau(1/\pi-1/2)},
  \qquad \tau\Bigl(\frac1\pi - \frac12\Bigr) = \frac{2/\pi - 1}{2s+1}.
\]

\emph{Caso $\pi \ge 2$.} Hölder em $\R^{2^{\lev}}$ dá
$\norm{\theta_{\lev}}_{2} \le 2^{\lev(1/2-1/\pi)}\norm{\theta_{\lev}}_{\pi}
\le C_{g}2^{-\lev s}$ (com $1/\pi = 0$ se $\pi = \infty$), e a segunda cota é
$r_{\lev} \le \norm{\theta_{\lev}}_{2}^{2} \le C_{g}^{2}2^{-2\lev s}$. Com
$2^{x(2s+1)} = C_{g}^{2}b/\eta$ e $h = 2^{x}\eta/b = C_{g}^{2}2^{-2xs}$, a
mesma conta dá a contribuição $\eta + \{2 + (1-2^{-2s})^{-1}\}h$, e
$h = (C_{g}^{2}b/\eta)^{1/(2s+1)}\eta/b = C_{g}^{\tau}(\eta/b)^{2s/(2s+1)}$,
porque $2/(2s+1) = \tau$.

Somar sobre os $pq$ blocos dá \eqref{eq:risco-pedacos}. Em $\pi = 2$,
$1 - 2^{1-a} = 1 - 2^{-2s}$ e as duas constantes coincidem. \qed

\subsection*{As duas consequências da Hipótese~\ref{ass:regime-blocos}}

\emph{\eqref{eq:ordem-lambda}.} $(\lamGone)^{2} = 4\max_{G}\lamzbar{G}^{2}
\le (8\sigma^{2}/n)\{2B_{X}^{2}C_{U}b_{\max} + 2\Lambda\log(|\Gcal_{n}|/\alpha)\}$
por $(x+y)^{2} \le 2x^{2}+2y^{2}$. Pela Proposição~\ref{prop:balanceada},
$b_{\max} = \max_{G}|G| \le 2\bn - 1 < 2\bn$ e $|\Gcal_{n}| \le pq(1 +
2^{J_{n}}/\bn) \le pq(1+n) \le 2pqn$, porque $2^{J_{n}} \le 2n$ e
$\bn \ge 2$; logo $\log(|\Gcal_{n}|/\alpha) \le \log(2pq/\alpha) + \log n
\le \log(2pq/\alpha) + \bn$, e
$(\lamGone)^{2} \le (16\sigma^{2}/n)\{(2B_{X}^{2}C_{U} + \Lambda)\bn +
\Lambda\log(2pq/\alpha)\} \le C_{\pen}\sigma^{2}\bn/n$, usando $\bn \ge 2$.

\emph{O evento $\Acal$.} Em $\Acal$, o Lema~\ref{lem:calib-blocos}(ii) dá
$\lamz{G} \le \lamzbar{G}$ para todo $G$, logo $\lamG \ge \pen_{w}$ e, pelo
mesmo lema, $\Ecal \subseteq \Tw$. Pela §1, $\lmin(\Gramhat) > \gamma/2$ e
$\lmaxx(\Gramhat) < \Lambda$; o Lema~\ref{lem:perfil-blocos}(iii) dá
$\gtil > \gamma/2$, e $\lmin(\Amat\T\Amat/n) \ge \lmin(\Gramhat) > 0$ dá o
posto. $\Prob(\smax^{2} \le 2B_{X}^{2}C_{U}) \to 1$ pelo Lema~6 de E1.5,
que pede $2^{J}\log d/n \to 0$, e $\Prob(\ESig) \to 1$ pela
Proposição~3(iv) de E1.4, que pede $pq2^{J}\log(pq2^{J})/n \to 0$; a segunda
condição implica a primeira.

\subsection*{Teorema~\ref{thm:taxa-blocos}}

\emph{(i).} Por \eqref{eq:Jn-blocos}, $n^{1/(2\seff+1)} \le 2^{J_{n}} <
2n^{1/(2\seff+1)} \le 2n$, logo a Hipótese~\ref{ass:regime-blocos} vale e
$pq2^{J_{n}}\log(pq2^{J_{n}})/n \asymp n^{1/(2\seff+1)-1}\log n =
n^{-2\seff/(2\seff+1)}\log n \to 0$; o resto é a segunda consequência acima.

\emph{(ii) e (iii).} Fixe $\epsilon > 0$ e tome $\alpha = \alpha' =
\alpha'' = \epsilon/4$. Intersecte $\Acal$, $\Ecal$,
$\{\normn{\bb}^{2} \le \mathcal{B}_{J_{n}}^{2}/\alpha'\}$ (Markov, porque
$\E\normn{\bb}^{2} = \norm{f - f_{J_{n}}}_{L_{2}(P)}^{2} \le
\mathcal{B}_{J_{n}}^{2}$ pelo Corolário~1 de E1.3) e
$\{\normn{\PA\bveps}^{2} \le \sigma^{2}p/(n\alpha'')\}$ (Markov, pelo
Teorema~\ref{thm:oraculo-blocos}(iii)); para $n$ grande, a interseção tem
probabilidade ao menos $1 - \epsilon$. Nela vale $\Tw$, e o
Teorema~\ref{thm:oraculo-blocos}(iii) com $\bar\bbeta = \betastar$ dá
\[
  \normn{\widehat f - f}^{2}
  \le \frac{192\,(\lamG)^{2}\mathcal{W}(\Gzero(\thetastar))}{\gtil}
  + 60\normn{\bb}^{2} + 3\normn{\PA\bveps}^{2}.
\]
Pelo Lema~\ref{lem:calib-blocos}(iii), por \eqref{eq:ordem-lambda} e pela
Proposição~\ref{prop:balanceada}(ii), com $\gtil > \gamma/2$,
\[
  \frac{192\,(\lamG)^{2}\mathcal{W}(\Gcal_{n})}{\gtil}
  \le \frac{384\,\bar\varrho^{2}(\lamGone)^{2}|\Gcal_{n}|}{\gamma}
  \le \frac{384\,\bar\varrho^{2}C_{\pen}\sigma^{2}\bn}{\gamma n}\,
  pq\Bigl(1 + \frac{2^{J_{n}}}{\bn}\Bigr)
  = \frac{384\,\bar\varrho^{2}C_{\pen}\sigma^{2}pq\,(\bn + 2^{J_{n}})}{\gamma n},
\]
que é $\asymp n^{-2\seff/(2\seff+1)}$, porque $2^{J_{n}} \asymp
n^{1/(2\seff+1)}$ domina $\bn \asymp \log n$. O viés é
$60\mathcal{B}_{J_{n}}^{2}/\alpha' \asymp 2^{-2J_{n}\seff} \le
n^{-2\seff/(2\seff+1)}$, pela cota inferior de $2^{J_{n}}$, e o último termo é
$O(1/n)$. A parcela $\bn/n$ e $\sigma^{2}p/n$ são de ordem estritamente
menor, porque $2\seff/(2\seff+1) < 1$. Logo
$\Prob(\normn{\widehat f - f}^{2} > M_{\epsilon}n^{-2\seff/(2\seff+1)}) <
\epsilon$ para uma constante $M_{\epsilon}$, que é (ii). \qed

\subsection*{Corolário~\ref{cor:componentes-blocos}}

Pela ortonormalidade das $\{\wav{\lev}{\tsl}\}_{\lev<J_{n}}$ em $L_{2}[0,1]$
(Corolário~3 de E1.5) e pelo Lema~1 de E1.3,
$\sum_{\cv\md}\lVert\widehat g_{\cv\md} - \gcomp{\cv}{\md}\rVert^{2} \le
2\norm{\thetahat - \thetastar}_{2}^{2} + 2pqC_{g}^{2}(1-2^{-2\seff})^{-1}
2^{-2J_{n}\seff}$. No evento da prova do Teorema~\ref{thm:taxa-blocos}, o
Teorema~\ref{thm:oraculo-blocos}(ii) com $\bar\bbeta = \betastar$ dá
$\norm{\thetahat - \thetastar}_{2}^{2} \le 64(\lamG)^{2}\mathcal{W}(\Gcal_{n})/\gtil^{2}
+ 16\normn{\bb}^{2}/\gtil$, com $\gtil > \gamma/2$: as duas parcelas têm a
ordem das do Teorema~\ref{thm:taxa-blocos}, com $\gamma^{-2}$ no lugar de
$\gamma^{-1}$. Para $\chat$, como na prova do Corolário~4 de E1.6:
por \eqref{eq:decomp-blocos} com $\bar\bbeta = \betastar$,
$\Amat(\chat - \bc) = \PA(\bb + \bveps) - \PA\Bmat\vbar$, logo
$\normn{\Amat(\chat - \bc)} \le \normn{\bb} + \normn{\PA\bveps} +
\sqrt{\Lambda}\norm{\vbar}_{2}$, e $\lmin(\Amat\T\Amat/n) > \gamma/2$ dá
$\norm{\chat - \bc}_{2}^{2} \le 2\normn{\Amat(\chat-\bc)}^{2}/\gamma$. A
passagem a $\fcoef{\cv}$ é a mesma de E1.6. \qed

\subsection*{Corolário~\ref{cor:adaptativa-blocos}}

A janela: $(2-\tau)/(4\seff) = (4s/(2s+1))/(4\seff) = s/((2s+1)\seff)$, e
$s/((2s+1)\seff) \ge 1/(2s+1) = \tau/2$ porque $\seff \le s$; ela é não vazia
se e só se $s/((2s+1)\seff) < 1$. Com $c < 1$, $2^{J_{n}} \le 2n^{c} \le 2n$, e
a Hipótese~\ref{ass:regime-blocos} vale; a condição de E1.4 é
$pq2^{J_{n}}\log(pq2^{J_{n}})/n \asymp n^{c-1}\log n \to 0$, logo
$\Prob(\Acal) \to 1$.

Fixe $\epsilon > 0$ e o mesmo evento da prova do
Teorema~\ref{thm:taxa-blocos}. Nele, o Corolário~\ref{cor:risco-ideal} com
$\Lambda' = \Lambda$ dá $\normn{\widehat f - f}^{2} \le
120\Lambda\Rid(\thetastar;\eta) + 120\normn{\bb}^{2} + 3\normn{\PA\bveps}^{2}$,
e pelo Lema~\ref{lem:calib-blocos}(iii), por \eqref{eq:ordem-lambda} e por
$\gtil > \gamma/2$,
\[
  \eta = \frac{1{,}6\max_{G}(\lamG w_{G})^{2}}{\Lambda\gtil}
  \;\le\; \frac{3{,}2\,\bar\varrho^{2}(\lamGone)^{2}}{\Lambda\gamma}
  \;\le\; \bar\eta_{n} := \kappa\,\frac{\bn}{n},
  \qquad \kappa = \frac{3{,}2\,\bar\varrho^{2}C_{\pen}\sigma^{2}}{\Lambda\gamma}.
\]
Como $\Rid$ é não decrescente em $\eta$, o Lema~\ref{lem:risco-pedacos} com
$b = \bn$ dá
\[
  \Rid(\thetastar;\eta) \le \Rid(\thetastar;\bar\eta_{n})
  \le pq\Bigl\{\AG C_{g}^{\tau}\kappa^{2s/(2s+1)}\,n^{-2s/(2s+1)}
  \bn^{(2/\pi-1)_{+}/(2s+1)} + \kappa\frac{\bn}{n}\Bigr\},
\]
e $\bn \asymp \log n$; a segunda parcela é de ordem menor porque
$2s/(2s+1) < 1$. O viés é $O(2^{-2J_{n}\seff}) = O(n^{-2c\seff})$ e
$2c\seff \ge 2s/(2s+1)$ pela janela; $\sigma^{2}p/n$ é de ordem menor. Para as
componentes, a segunda cota de \eqref{eq:risco-ideal} e o Lema~1 de E1.3,
como no Corolário~\ref{cor:componentes-blocos}. Os itens: (i) e (ii) são o
expoente $(2/\pi-1)_{+}$; (iii) é
$2s/(2s+1) - (2/\pi-1)_{+}/(2s+1) = 2\seff/(2s+1)$, porque
$2s - 2/\pi + 1 = 2(s - 1/\pi + 1/2) = 2\seff$ em $\pi < 2$ e $\seff = s$ em
$\pi \ge 2$; (iv) é $\kappa^{2s/(2s+1)} \propto \bar\varrho^{\,4s/(2s+1)}$. \qed

\subsection*{Corolário~\ref{cor:lenta-blocos}}

\emph{A cota não assintótica \eqref{eq:lenta-nao-assint}.} No evento $\Tw$, o
Teorema~\ref{thm:oraculo-blocos}(i) com $\bar\bbeta = \betastar$ (logo
$\bbar = \bb$ e $\vbar = v := \thetahat - \thetastar$) dá
$\normn{\Btil v}^{2} \le 6\pen\nGw{\thetastar} + 4\normn{\bb}^{2}$, e a
primeira parte do Teorema~\ref{thm:oraculo-blocos}(iii),
$\normn{\widehat f - f} \le \normn{\Btil v} + 2\normn{\bb} +
\normn{\PA\bveps}$, com $(x+y+z)^{2} \le 3(x^{2}+y^{2}+z^{2})$, dá
$\normn{\widehat f - f}^{2} \le 3(6\pen\nGw{\thetastar} + 4\normn{\bb}^{2} +
4\normn{\bb}^{2} + \normn{\PA\bveps}^{2})$, que é
\eqref{eq:lenta-nao-assint}. Os Passos~1 a~3 e~6 da prova do teorema, os
únicos usados, não envolvem $\Gramhat$.

\emph{O $\pen$ da teoria e a forma $O_{p}$.} Seja $\Acal^{+} = \{\smax^{2} \le
2B_{X}^{2}C_{U}\} \cap \{\lmaxx(\Gramhat) \le \Lambda\}$, que contém $\Acal$.
Nele o Lema~\ref{lem:calib-blocos}(ii) dá $\lamz{G} \le \lamzbar{G}$ para todo
$G$, logo $\lamG \ge \pen_{w}$ e $\Ecal \subseteq \Tw$.
$\Prob(\smax^{2} \le 2B_{X}^{2}C_{U}) \to 1$ pelo Lema~6 de E1.5, que pede
$2^{J_{n}}\log d_{n}/n \to 0$, e $\Prob(\lmaxx(\Gramhat) \le \Lambda) \ge
\Prob(\opnorm{\Gramhat - \Gram} < \Lambda - \kappa_{2}C_{U}) \to 1$ pela cauda da
Proposição~3(iv) de E1.4 e por Weyl, que pedem
$pq2^{J_{n}}\log(pq2^{J_{n}})/n \to 0$; esta implica aquela. Fixe
$\epsilon > 0$, tome $\alpha = \alpha' = \alpha'' = \epsilon/4$ e intersecte
$\Acal^{+}$, $\Ecal$, $\{\normn{\bb}^{2} \le \mathcal{B}_{J_{n}}^{2}/\alpha'\}$ e
$\{\normn{\PA\bveps}^{2} \le \sigma^{2}p/(n\alpha'')\}$ (Markov, como na prova
do Teorema~\ref{thm:taxa-blocos}); para $n$ grande a interseção tem
probabilidade ao menos $1 - \epsilon$, e nela
$\normn{\widehat f - f}^{2} \le 18\lamG\nGw{\thetastar} +
24\mathcal{B}_{J_{n}}^{2}/\alpha' + 3\sigma^{2}p/(n\alpha'')$, que é a primeira
parte de \eqref{eq:lenta-blocos}. A cadeia: pelo
Lema~\ref{lem:calib-blocos}(iii), $\lamG w_{G} \le \varrho\lamGone$ para todo
$G$, e somar com os pesos $\norm{\thetastar_{G}}_{2}$ dá
$\lamG\nGw{\thetastar} \le \varrho\lamGone\nGone{\thetastar}$, com
$\varrho \le \bar\varrho$; e $\norm{x}_{2} \le \norm{x}_{1}$ em cada pedaço dá
$\nGone{\thetastar} \le \norm{\thetastar}_{1}$. A ordem de $\lamGone$ é
\eqref{eq:ordem-lambda}.

\emph{A contagem \eqref{eq:contagem-lenta}.} Fixe um bloco e escreva
$\theta_{\lev}$ para os seus coeficientes no nível $\lev$, como na prova do
Lema~\ref{lem:risco-pedacos}. Pela Hipótese de Besov,
$\norm{\theta_{\lev}}_{2} \le 2^{\lev(1/2-1/\pi)_{+}}\norm{\theta_{\lev}}_{\pi}
\le C_{g}2^{-\lev\seff}$ (Hölder em $\R^{2^{\lev}}$ se $\pi \ge 2$;
$\norm{\cdot}_{2} \le \norm{\cdot}_{\pi}$ se $\pi \le 2$). O pedaço grosso
contribui $\norm{\theta_{\mathrm{grosso}}}_{2} \le
\sum_{\lev\le\jst}\norm{\theta_{\lev}}_{2} \le C_{g}(1 - 2^{-\seff})^{-1}$. Num
nível fino $\lev$, os $k_{\lev} \le 2^{\lev}/b$ pedaços $G \in \Gcal_{\lev}$
particionam o nível. Se $\pi \ge 2$, por Cauchy--Schwarz,
\[
  \sum_{G\in\Gcal_{\lev}}\norm{\theta_{G}}_{2}
  \le k_{\lev}^{1/2}\norm{\theta_{\lev}}_{2}
  \le \Bigl(\frac{2^{\lev}}{b}\Bigr)^{1/2}C_{g}2^{-\lev s}
  = C_{g}\,b^{-1/2}\,2^{\lev(1/2-s)};
\]
se $\pi \le 2$, por $\norm{\theta_{G}}_{2} \le \norm{\theta_{G}}_{\pi}$ e Hölder
entre os pedaços,
\[
  \begin{aligned}
  \sum_{G\in\Gcal_{\lev}}\norm{\theta_{G}}_{2}
  &\le k_{\lev}^{1-1/\pi}\Bigl(\sum_{G\in\Gcal_{\lev}}\norm{\theta_{G}}_{\pi}^{\pi}\Bigr)^{1/\pi}
  = k_{\lev}^{1-1/\pi}\norm{\theta_{\lev}}_{\pi}\\
  &\le \Bigl(\frac{2^{\lev}}{b}\Bigr)^{1-1/\pi}C_{g}2^{-\lev(s+1/2-1/\pi)}
  = C_{g}\,b^{-(1-1/\pi)}\,2^{\lev(1/2-s)}.
  \end{aligned}
\]
As duas são $C_{g}b^{-\vartheta_{\pi}}2^{\lev(1/2-s)}$, e são atingidas: com a
energia do nível espalhada igualmente, em $\pi \ge 2$, e com uma espiga por
pedaço, em $\pi \le 2$. Somar os níveis finos $\jst < \lev < J_{n}$ e os $pq$
blocos, com $b = \bn$, dá \eqref{eq:contagem-lenta}.

\emph{(i).} Com $s > 1/2$ e $2^{\jst+1} \ge b$,
$\sum_{\lev>\jst}2^{\lev(1/2-s)} \le 2^{(\jst+1)(1/2-s)}/(1 - 2^{1/2-s}) \le
\bn^{1/2-s}/(1 - 2^{1/2-s})$, e \eqref{eq:contagem-lenta} dá
$\nGone{\thetastar} \le pqC_{g}\{(1 - 2^{-\seff})^{-1} +
\bn^{1/2-s-\vartheta_{\pi}}(1 - 2^{1/2-s})^{-1}\} = O(1)$; logo
$\lamG\nGw{\thetastar} = O(\sqrt{\bn/n}) = O((\log n/n)^{1/2})$. Com
$J_{n} \ge \log_{2}n/(4\seff)$, $\mathcal{B}_{J_{n}}^{2} \asymp
2^{-2J_{n}\seff} \le n^{-1/2}$, e $p/n$ é de ordem menor. Com $\seff > 1/4$, a
escolha $J_{n} = \lceil\log_{2}n/(4\seff)\rceil$ tem $2^{J_{n}} <
2n^{1/(4\seff)}$ e satisfaz a condição de dimensão.

\emph{(ii).} Com $s < 1/2$, $\sum_{\jst<\lev<J_{n}}2^{\lev(1/2-s)} \le
2^{J_{n}(1/2-s)}/(2^{1/2-s} - 1)$, logo $\nGone{\thetastar} =
O(1 + \bn^{-\vartheta_{\pi}}2^{J_{n}(1/2-s)})$ e, como $1/2 - \vartheta_{\pi} =
(1/\pi - 1/2)_{+}$,
\[
  \lamG\nGw{\thetastar} \;\lesssim\; \sqrt{\frac{\bn}{n}}
  + n^{-1/2}\,\bn^{(1/\pi-1/2)_{+}}\,2^{J_{n}(1/2-s)} .
\]
Igualar a segunda parcela ao viés $2^{-2J_{n}\seff}$ dá
$2^{J_{n}(1/2-s+2\seff)} \asymp n^{1/2}\bn^{-(1/\pi-1/2)_{+}}$, que é a escolha
do enunciado, com valor comum $2^{-2J_{n}\seff} \asymp
(\bn^{(2/\pi-1)_{+}}/n)^{2\seff/(1-2s+4\seff)}$. A parcela $\sqrt{\bn/n}$ e
$p/n$ são de ordem menor, porque $2\seff/(1-2s+4\seff) < 1/2$ equivale a
$s < 1/2$. A condição de dimensão vale porque $\seff > s/2$ dá
$1/(1-2s+4\seff) < 1$; e $2^{J_{n}} \ge \bn$ para $n$ grande. Em $\pi \ge 2$,
$\seff = s$, $(2/\pi-1)_{+} = 0$ e $1 - 2s + 4s = 2s + 1$.

\emph{Os dois eventos acima do teto (para (iii) e (iv)).} Seja $\Acal^{++} =
\{\smax^{2} \le 2B_{X}^{2}C_{U}\mu^{\Gcal}_{J}\} \cap \{\lmaxx(\Gramhat) \le
\Lambda\mu^{\Gcal}_{J}\}$. A conta da prova da Proposição~4(iii) de E1.6 dá
$\Prob(\smax^{2} > 2B_{X}^{2}C_{U}(1 + c_{\psi}2^{J}\log d/n)) \le 1/d$, e
$c_{\psi}2^{J}\log d \le c_{\mu}2^{J}\log(2D)$. Para a segunda, a cauda da
Proposição~3(iv) de E1.4, $\Prob(\opnorm{\Gramhat - \Gram} \ge t) \le
2D\exp\{-nt^{2}/(2R_{J}(\kappa_{2}C_{U} + 2t/3))\}$, vale para todo $J$ e
todo $t > 0$; com $x = 2\log(2D)$ e $t = (2xR_{J}\kappa_{2}C_{U}/n)^{1/2} +
(4/3)xR_{J}/n$, vale $nt^{2} \ge 2xR_{J}\kappa_{2}C_{U} + (4/3)xR_{J}t$, e a
probabilidade é no máximo $2De^{-x} = 1/(2D)$; e $(2ab)^{1/2} \le \xi +
ab/(2\xi)$ com $a = \kappa_{2}C_{U}$ e $b = xR_{J}/n$ dá $t \le \xi +
(\kappa_{2}C_{U}/(2\xi) + 4/3)xR_{J}/n$. Por Weyl e
$\lmaxx(\Gram) \le \kappa_{2}C_{U}$, fora desse evento
$\lmaxx(\Gramhat) < \Lambda + (\kappa_{2}C_{U}/(2\xi) + 4/3)
2\log(2D)R_{J}/n \le \Lambda\mu^{\Gcal}_{J}$, porque $R_{J} =
pB_{X}^{2}(1 + qC_{\psi}2^{J}) \le pB_{X}^{2}(1+qC_{\psi})2^{J}$. Logo
$\Prob(\Acal^{++}) \ge 1 - 1/d - 1/(2D) \to 1$, sem condição sobre $J_{n}$.
No evento $\Acal^{++}$, o Lema~\ref{lem:calib-blocos}(ii) dá $\lamz{G} \le
(\mu^{\Gcal}_{J})^{1/2}\lamzbar{G}$ para todo $G$, logo $\lamGp \ge
\pen_{w}$ e $\Ecal \subseteq \Tw$, e o passo da forma $O_{p}$ acima segue com
$\lamGp$ no lugar de $\lamG$ e $\Acal^{++}$ no lugar de $\Acal^{+}$; a cadeia
dá $\lamGp\nGw{\thetastar} \le \bar\varrho(\mu^{\Gcal}_{J_{n}})^{1/2}
\lamGone\nGone{\thetastar}$. Com $2^{J_{n}} \le n^{2}$, $|\Gcal_{n}| \le pq(1 +
n^{2}/\bn) \le pqn^{2}$ e $\log(|\Gcal_{n}|/\alpha) \le \log(2pq/\alpha) +
2\bn$, e a conta de \eqref{eq:ordem-lambda} dá $(\lamGone)^{2} \le
(16\sigma^{2}/n)\{(2B_{X}^{2}C_{U} + 2\Lambda)\bn + \Lambda\log(2pq/\alpha)\}
\le (C_{\pen} + 16\Lambda)\sigma^{2}\bn/n$. Como $\log(2D_{n}) \asymp \log n
\asymp \bn$,
\begin{equation}
  \label{eq:lam-plus-blocos}
  (\mu^{\Gcal}_{J_{n}})^{1/2}\lamGone \;\asymp\;
  \Bigl(\frac{\log n}{n}\Bigr)^{1/2}\Bigl(1 + \frac{2^{J_{n}}\log n}{n}\Bigr)^{1/2},
  \qquad\text{e}\qquad
  \asymp 2^{J_{n}/2}\,\frac{\log n}{n}\quad\text{se } 2^{J_{n}} \gtrsim n/\log n .
\end{equation}

\emph{(iii)(a).} Com $s > 1/2$, $\nGone{\thetastar} = O(1)$ como em (i), e a
conta é a da Proposição~4(iii)(a) de E1.6 com
\eqref{eq:lam-plus-blocos} no lugar da ordem de $\pen_{n}^{+}$ de lá: com
$2^{J_{n}} \asymp (n/\log n)^{2/(1+4\seff)} \gtrsim n/\log n$, o termo de
estimação é $(\log n/n)(n/\log n)^{1/(1+4\seff)} =
(\log n/n)^{4\seff/(1+4\seff)}$, a ordem do viés.

\emph{(iii)(b).} Com $s < 1/2$ e $\pi < 2$, $\vartheta_{\pi} = 1 - 1/\pi$ e
$\nGone{\thetastar} = O(1 + \bn^{-\vartheta_{\pi}}2^{J_{n}(1/2-s)})$ como em
(ii). Com $2^{J_{n}} \asymp (n/\bn^{1/\pi})^{1/(1-s+2\seff)}$, que passa de
$n/\bn$ para $n$ grande (o expoente é pelo menos $1$ e $1/\pi \le 1$),
\eqref{eq:lam-plus-blocos} dá o termo de estimação
\[
  2^{J_{n}/2}\frac{\log n}{n}\bigl(1 + \bn^{-(1-1/\pi)}2^{J_{n}(1/2-s)}\bigr)
  \;\asymp\; 2^{J_{n}/2}\frac{\log n}{n} +
  2^{J_{n}(1-s)}\frac{(\log n)^{1/\pi}}{n},
\]
cuja segunda parcela domina e, na escolha do enunciado, é
\[
  (\bn^{1/\pi}/n)^{1 - (1-s)/(1-s+2\seff)} =
  (\bn^{1/\pi}/n)^{2\seff/(1-s+2\seff)},
\]
a ordem do viés $2^{-2J_{n}\seff}$.

\emph{(iv).} Com $s = 1/2$, \eqref{eq:contagem-lenta} tem
$\sum_{\jst<\lev<J_{n}}2^{0} \le J_{n}$, o que dá a cota de
$\nGone{\thetastar}$ do enunciado, e $J_{n} \asymp \log n \asymp \bn$ nas
duas escolhas, de modo que $\nGone{\thetastar} \asymp 1 +
(\log n)^{1-\vartheta_{\pi}} \asymp (\log n)^{1-\vartheta_{\pi}}$
($\vartheta_{\pi} \le 1/2$). Se $\seff \ge 1/4$, $2^{J_{n}} \asymp
(n/\bn)^{1/(4\seff)} \lesssim n/\log n$, $\mu^{\Gcal}_{J_{n}}$ é limitado, o
termo de estimação é $(\log n/n)^{1/2}(\log n)^{1-\vartheta_{\pi}}$ e o viés,
$(\bn/n)^{1/2}$, é de ordem menor. Se $\seff < 1/4$, $\pi < 4/3$ e
$\vartheta_{\pi} = 1 - 1/\pi$; com $2^{J_{n}} \asymp
(n/\bn^{1+1/\pi})^{2/(1+4\seff)}$, que passa de $n/\log n$, o termo de
estimação é $2^{J_{n}/2}(\log n)^{1+1/\pi}/n \asymp
((\log n)^{1+1/\pi}/n)^{4\seff/(1+4\seff)}$, a ordem do viés. \qed

\subsection*{Corolário~\ref{cor:truncada-blocos}}

\emph{(i).} O Teorema~\ref{thm:oraculo-blocos}(i) e (iii) com $\bar\bbeta =
(\bc,\thbar)$ dão, como na prova de \eqref{eq:lenta-nao-assint},
$\normn{\widehat f - f}^{2} \le 18\pen\nGw{\thbar} + 24\normn{\bbar}^{2} +
3\normn{\PA\bveps}^{2}$, e $\bbar = \bb + \Bmat(\thetastar - \thbar)$ dá
$\normn{\bbar}^{2} \le 2\normn{\bb}^{2} + 2\normn{\Bmat(\thetastar -
\thbar)}^{2}$. A cota em esperança é a da Proposição~8(i) de E1.6, porque
$\thetastar - \thbar$ é determinístico quando $\pen$ e $w$ o são. Nos pedaços
com $\norm{\thetastar_{G}}_{2} > \pen w_{G}$, $\pen w_{G}\norm{\thetastar_{G}}_{2}$
é o mínimo de \eqref{eq:RG}; nos outros, $\norm{\thetastar_{G}}_{2}^{2}$ é; somar
dá a cota por $\max(18,48\kappa_{2}C_{U})\mathcal{R}_{\Gcal,w}$. Pelo
Lema~\ref{lem:calib-blocos}(iii), $\lamG w_{G} \le \varrho\lamGone$, e
$\min(\varrho u, v) \le \varrho\min(u,v)$ para $\varrho \ge 1$ dá
$\mathcal{R}_{\Gcal,w}(\thetastar;\lamG) \le
\varrho\,\mathcal{R}_{\Gcal,1}(\thetastar;\lamGone)$; com $\lamGp$, o mesmo com
$(\mu^{\Gcal}_{J})^{1/2}\lamGone$. A forma $O_{p}$ é a da prova do
Corolário~\ref{cor:lenta-blocos}: os eventos $\Acal^{+}$ (abaixo do teto) ou
$\Acal^{++}$ (em todo $J$), $\Ecal \subseteq \Tw$, e Markov em
$\normn{\bb}^{2}$, $\normn{\Bmat(\thetastar - \thbar)}^{2}$ e
$\normn{\PA\bveps}^{2}$.

\emph{(ii), $s < 1/2$.} Fixe um bloco. O pedaço grosso contribui no máximo
$t\norm{\theta_{\mathrm{grosso}}}_{2} \le tC_{g}(1-2^{-\seff})^{-1}$, pela
prova de \eqref{eq:contagem-lenta}. Num nível fino $\lev$, com $\Gcal_{\lev}$
os seus pedaços e $v_{\lev} = \sum_{G\in\Gcal_{\lev}}\min(t\norm{\theta_{G}}_{2},
\norm{\theta_{G}}_{2}^{2})$, há duas cotas. A primeira é a de grupos:
$v_{\lev} \le t\sum_{G}\norm{\theta_{G}}_{2} \le tC_{g}b^{-\vartheta_{\pi}}
2^{\lev(1/2-s)} =: A_{\lev}$, pela prova de \eqref{eq:contagem-lenta}. A
segunda é a cauda: se $\pi \le 2$, $\min(ty,y^{2}) \le t^{2-\pi}y^{\pi}$ e
$\norm{\theta_{G}}_{2} \le \norm{\theta_{G}}_{\pi}$ dão $v_{\lev} \le
t^{2-\pi}\sum_{G}\norm{\theta_{G}}_{\pi}^{\pi} =
t^{2-\pi}\norm{\theta_{\lev}}_{\pi}^{\pi} \le t^{2-\pi}C_{g}^{\pi}2^{-\lev\pi\seff}$;
se $\pi \ge 2$, $v_{\lev} \le \sum_{G}\norm{\theta_{G}}_{2}^{2} =
\norm{\theta_{\lev}}_{2}^{2} \le C_{g}^{2}2^{-2\lev s}$. É a mesma cota $B_{\lev}$
da prova da Proposição~8(ii) de E1.6; só $A_{\lev}$ ganhou o fator
$b^{-\vartheta_{\pi}}$. Com
\[
  2^{x} = \bigl\{(C_{g}/t)\,b^{1-\vartheta_{\pi}}\bigr\}^{\tau},
  \qquad
  h = C_{g}^{\tau}\,t^{2-\tau}\,b^{\tau(1-\vartheta_{\pi})-1},
\]
vale $A_{\lev} = h2^{(\lev-x)(1/2-s)}$ e $B_{\lev} = h2^{-(\lev-x)\min(\pi,2)\seff}$:
em $\lev = x$,
\[
  A_{x} = C_{g}^{\tau}t^{2-\tau}b^{-\vartheta_{\pi} + (1-\vartheta_{\pi})(\tau-1)},
\]
porque $\tau(1/2-s) = \tau - 1$, e o expoente de
$b$ é $\tau(1-\vartheta_{\pi}) - 1$; se $\pi \le 2$, $1 - \vartheta_{\pi} =
1/\pi$ e
\[
  B_{x} = t^{2-\pi}C_{g}^{\pi}\{(C_{g}/t)b^{1/\pi}\}^{-(\pi-\tau)} =
  C_{g}^{\tau}t^{2-\tau}b^{\tau/\pi-1};
\]
se $\pi \ge 2$, $1 - \vartheta_{\pi} =
1/2$ e $B_{x} = C_{g}^{2}\{(C_{g}/t)b^{1/2}\}^{-(2-\tau)} =
C_{g}^{\tau}t^{2-\tau}b^{\tau/2-1}$. Em todo $\pi$, $\tau(1-\vartheta_{\pi}) - 1 =
((2/\pi-1)_{+}-2s)/(2s+1)$, e $2 - \tau = 4s/(2s+1)$. Somar $A_{\lev}$ nos níveis
finos $\lev \le x$ e $B_{\lev}$ nos $\lev > x$, como na Proposição~8 de E1.6,
dá no máximo $h\{(1-2^{s-1/2})^{-1} + (1-2^{-\min(\pi,2)\seff})^{-1}\}$ por
bloco. Truncar em $\lev < J$ só tira parcelas, e somar os $pq$ blocos dá
\eqref{eq:contagem-truncada-blocos}.

\emph{(ii), $s \ge 1/2$.} A primeira afirmação é $\min(tu,u^{2}) \le tu$. Em
$s > 1/2$, o bloco com $\theta^{*}_{\cv\md,00} = C_{g}$ e o resto nulo tem
pedaço grosso de norma $C_{g}$, que contribui $\min(tC_{g},C_{g}^{2}) = tC_{g}$.
Em $s = 1/2$, num nível fino com $2^{\lev} \le C_{g}b^{1-\vartheta_{\pi}}/t$:
se $\pi \le 2$, ponha uma espiga de altura $x_{\lev} = C_{g}2^{-\lev(1-1/\pi)}
k_{\lev}^{-1/\pi}$ em cada um dos $k_{\lev}$ pedaços, o que satura a Hipótese de
Besov no nível; como $k_{\lev} \le 2^{\lev}/b$, $x_{\lev} \ge
C_{g}2^{-\lev}b^{1/\pi} \ge t$, e o nível contribui $k_{\lev}tx_{\lev} =
tC_{g}2^{-\lev(1-1/\pi)}k_{\lev}^{1-1/\pi} \ge tC_{g}(2b)^{-(1-1/\pi)}$, porque
$k_{\lev} \ge 2^{\lev}/(2b)$; se $\pi \ge 2$, espalhe $C_{g}2^{-\lev}$ em todas
as coordenadas: cada pedaço tem norma $|G|^{1/2}C_{g}2^{-\lev} \ge
b^{1/2}C_{g}2^{-\lev} \ge t$, e o nível contribui
$tC_{g}2^{-\lev}\sum_{G}|G|^{1/2} \ge tC_{g}2^{-\lev}k_{\lev}b^{1/2} \ge
tC_{g}b^{-1/2}/2$. Nos dois casos, pelo menos $tC_{g}b^{-\vartheta_{\pi}}/2$.

\emph{(iii).} Com $\pen = \lamG$ e $(\lamGone)^{2} \le C_{\pen}\sigma^{2}\bn/n$
\eqref{eq:ordem-lambda}, \eqref{eq:contagem-truncada-blocos} com $b = \bn$
dá $\mathcal{R}_{\Gcal,1}(\thetastar;\lamGone) \lesssim (\bn/n)^{1/2} +
(\bn/n)^{2s/(2s+1)}\bn^{((2/\pi-1)_{+}-2s)/(2s+1)} = (\bn/n)^{1/2} +
n^{-2s/(2s+1)}\bn^{(2/\pi-1)_{+}/(2s+1)}$, e a primeira parcela é de ordem
menor porque $2s/(2s+1) < 1/2$. O viés e a condição de dimensão são os da
prova do Corolário~\ref{cor:adaptativa-blocos}: $c \ge s/((2s+1)\seff)$ dá
$\mathcal{B}_{J_{n}}^{2} \lesssim n^{-2s/(2s+1)}$, e $c < 1$ dá
$pq2^{J_{n}}\log(pq2^{J_{n}})/n \to 0$ e $2^{J_{n}} \le 2n$.

\emph{(iv).} Escreva $\eta = 2s/(s+\seff(2s+1)) \in [1,2)$ e $\beta =
(2s+2/\pi-1)/(2s+1)$, que está em $(0,\, 4s/(2s+1))$ porque $\pi < 2$ e
$2\seff = 2s+1-2/\pi > 0$. A escolha do enunciado é $2^{J_{n}} \asymp
n^{\eta}\bn^{-\eta\beta(2s+1)/(4s)}$, que passa de $n/\bn$ porque
$\beta(2s+1)/(4s) < 1$ e tem $2^{J_{n}} \le n^{2}$. Por
\eqref{eq:lam-plus-blocos}, $t = (\mu^{\Gcal}_{J_{n}})^{1/2}\lamGone \asymp
2^{J_{n}/2}\bn/n$, e \eqref{eq:contagem-truncada-blocos} dá
\[
  t^{4s/(2s+1)}\bn^{\beta - 4s/(2s+1)}
  \;\asymp\; 2^{J_{n}\cdot 2s/(2s+1)}\,n^{-4s/(2s+1)}\,\bn^{\beta},
\]
que, igualado ao viés $2^{-2J_{n}\seff}$, dá $2^{J_{n}(2s/(2s+1)+2\seff)} =
n^{4s/(2s+1)}\bn^{-\beta}$: é a escolha do enunciado, porque
$(2s/(2s+1) + 2\seff)(2s+1) = 2\{s + \seff(2s+1)\}$. O valor comum é
$2^{-2J_{n}\seff} = (\bn^{\beta}/n^{4s/(2s+1)})^{2\seff/(2s/(2s+1)+2\seff)} =
n^{-e}\bn^{e\beta(2s+1)/(4s)}$, e $\beta(2s+1)/(4s) = (2s+2/\pi-1)/(4s)$. O
termo do pedaço grosso, $t \asymp 2^{J_{n}/2}\bn/n$, é polinomialmente menor,
porque $t \to 0$ e $4s/(2s+1) < 1$. \qed

\section{O que foi transposto e o que é novo}

\begin{itemize}[nosep, leftmargin=*]
  \item \textbf{Transposto de E1.5 e E1.6} (e, por eles, de Bühlmann e van de
    Geer, 2011, Seção~6.2, e da organização em passos do WALL teórico): os
    seis passos do Teorema~\ref{thm:oraculo-blocos}, a perfilagem, o viés por
    Young com os dois casos, a leitura com comparador arbitrário e o
    risco ideal como comparador truncado.
  \item \textbf{Transposto de Lounici et al.\ (2011):} a norma dual da norma
    de grupos e a decomponibilidade; a observação de que grupos com pelo
    menos $\log M$ coordenadas evitam o logaritmo que o LASSO paga
    (Teorema~7.1 deles, a cota inferior do LASSO no modelo multitarefa).
  \item \textbf{Transposto de Hsu, Kakade e Zhang (2012):} o Teorema~2.1, no
    Lema~\ref{lem:calib-blocos}(i).
  \item \textbf{Transposto de Klopp e Pensky (2015):} os pedaços de tamanho
    $\log n$, a forma do risco ideal por pedaços (Lema~4 do suplemento) e a
    taxa; e, de Cai (1999), a origem dos blocos de tamanho $\log n$ em
    limiarização de wavelets.
  \item \textbf{Novo:} a calibração em blocos com o desenho dependente e a
    forma fechada controlada pelos eventos de E1.4 e E1.5; o preço dos pesos
    como um fator $\varrho$ no termo de estimação
    (Lema~\ref{lem:calib-blocos}(iii)), que cobre os pesos $\sqrt{|G|}$ na
    forma balanceada; o oráculo sem cone para grupos com os níveis livres; o
    Lema~\ref{lem:risco-pedacos} sob a hipótese de Besov por nível, mais
    fraca que a (A3), com o pedaço grosso custando um $\eta$; as taxas com
    $p$ fixo, sem (A6), sem norma dual e com $\seff > s/(2s+1)$ no lugar de
    $r^{*} > (2\varsigma)^{-1}$; e a Observação~\ref{rem:branca}, que cobre
    o estimador exato do \texttt{grpreg}.
  \item \textbf{Novo em E1.13:} a contagem \eqref{eq:contagem-lenta} de
    $\nGone{\thetastar}$ sob Besov, justa em $b$, e com ela a taxa lenta sem
    logaritmo em $s < 1/2$ e $\pi \ge 2$ (Corolário~\ref{cor:lenta-blocos});
    a leitura da Proposição~4 de E1.6 pelo Teorema~\ref{thm:oraculo-blocos}(i)
    é transposta.
  \item \textbf{Novo em E1.14:} os itens (iii) e (iv) do
    Corolário~\ref{cor:lenta-blocos}, em todo $\seff > 0$ e na fronteira
    $s = 1/2$, com o nível acima do teto e $\lamG$ inflado por
    $(\mu^{\Gcal}_{J})^{1/2}$, que cobre ao mesmo tempo $\smax$ e
    $\lmaxx(\Gramhat)$ acima do teto \eqref{eq:mu-blocos}; o ganho em logaritmo
    sobre o LASSO em (iii)(b) e em (iv).
  \item \textbf{Novo em E1.15:} o Corolário~\ref{cor:truncada-blocos} e a
    contagem \eqref{eq:contagem-truncada-blocos} de
    $\mathcal{R}_{\Gcal,1}(\thetastar;t)$ sob Besov, justa em $t$ e em $b$, que
    junta a contagem do Corolário~\ref{cor:lenta-blocos} com o corte do
    Lema~\ref{lem:risco-pedacos}; com ela a taxa do
    Corolário~\ref{cor:adaptativa-blocos} sai, em $s < 1/2$, sem condição de
    desenho. O comparador truncado por pedaço é o do
    Corolário~\ref{cor:risco-ideal}, lido no Teorema~\ref{thm:oraculo-blocos}(i)
    em vez de (ii).
\end{itemize}

\section{O que isto não cobre}

\begin{itemize}[nosep, leftmargin=*]
  \item \textbf{$\pen$ e $t$ por validação cruzada.} Como em E1.5 e E1.6, a
    teoria é para $\pen$ determinístico (ou o $\pen_{w}$ calibrado em
    $(\bX,\bU)$). O WAFC escolhe $\pen$ por validação cruzada e o limiar $t$
    pela regra \texttt{cv1se} (D45); nenhum dos dois tem teorema aqui.
  \item \textbf{Pesos que dependem das respostas.} O argumento pede $w$ e, na
    Observação~\ref{rem:branca}, $Q_{G}$ mensuráveis em $(\bX,\bU)$, e
    $\varrho$ limitado. Pesos estimados a partir de um primeiro ajuste, que
    dependem de $\bY$, não são cobertos.
  \item \textbf{Pedaços que atravessam níveis} (a norma (3.1) literal de Klopp
    e Pensky). O Lema~\ref{lem:risco-pedacos} usa pedaços finos contidos num
    nível. Com pedaços atravessando níveis cada um toca no máximo dois níveis e
    a contagem muda por constante, mas isso não foi escrito.
  \item \textbf{$\bn$ diferente de $\lceil\log n\rceil$.} A
    Hipótese~\ref{ass:regime-blocos} fixa $\bn = \lceil\log n\rceil$. Com
    $\bn$ de ordem menor que $\log|\Gcal|$ volta parte do logaritmo; com
    $\bn$ maior, o fator $\bn^{(2/\pi-1)/(2s+1)}$ cresce. A constante ótima
    não foi estudada.
  \item \textbf{A cota inferior para $q \ge 2$} (E1.9). A
    Observação~\ref{rem:inferior} só empresta a de Klopp e Pensky para
    $q = 1$ e $\bX\perp\bU$.
  \item \textbf{$\sigma$ desconhecido, $p$ crescente, a base do intervalo e a
    norma populacional}, pelas mesmas razões de E1.5 e E1.6. Com $p$
    crescente, $\log|\Gcal|$ deixa de ser $O(\bn)$ e a (A6) de Klopp e Pensky
    volta a ser o regime natural.
  \item \textbf{Seleção pelo estimador.} O block LASSO zera pedaços, não
    blocos $(\cv,\md)$, e zerar todos os pedaços de um bloco nulo pediria uma
    condição de irrepresentabilidade (E1.7a), que os pedaços não mudam. A
    seleção de estrutura continua sendo estimação seguida de limiar, no
    adendo de E1.12 de \texttt{06-selecao-limiar.tex}.
  \item \textbf{Os níveis grossos livres} (\texttt{klopp.freecoarse}, §12 do
    \texttt{08a}): cobertos pela mesma prova com $\Amat$ aumentado pelas
    colunas dos níveis $\lev \le \jst$ e o termo $\sigma^{2}p_{0}/n$,
    $p_{0} = p + pq(2^{\jst+1}-1)$, que é $O(\log n/n)$; não enunciado aqui,
    porque D44 não adotou essa forma.
  \item \textbf{A taxa lenta além do comparador truncado.} O
    Corolário~\ref{cor:truncada-blocos} esgota a rota do
    Teorema~\ref{thm:oraculo-blocos}(i) em $s \ge 1/2$ (o pedaço grosso paga
    $\lamGone$ ou $O(1)$) e é justo em $t$ e $b$ em $s < 1/2$. Não dá as
    componentes, que pedem $\lmin(\Gramhat)$. A variante branca da
    Observação~\ref{rem:branca} pede o truncamento em $\pen w_{G}\Lambda^{1/2}$
    (determinístico) no evento $q_{\max} \le \Lambda$; a conta é a mesma, mas não
    foi escrita nem conferida.
  \item \textbf{$\smax$ e $\lmaxx(\Gramhat)$ muito acima do teto.} As cotas de
    \eqref{eq:mu-blocos} têm $2^{J}\log(2D)/n$, de Bernstein; quando $2^{J}/n$
    cresce como potência de $n$, as duas são da ordem de $2^{J}/n$, o que
    tiraria parte do logaritmo de (iii)(a) e de (iv) com $\seff < 1/4$. A
    constante $c_{\mu}$ (perto de $150$ no desenho da conferência) vem do
    $\xi$ pequeno de E1.6 e não foi otimizada.
  \item \textbf{A calibração sem $\lmaxx(\Gramhat)$.} Nos pedaços finos, as
    wavelets de um mesmo nível só se sobrepõem em translações vizinhas, e
    $\opnorm{\Psitil_{G}}$ fica sob uma constante da família vezes
    $\smax^{2}$, sem $\Lambda$; o pedaço grosso pediria uma concentração na
    dimensão $2^{\jst+1} - 1 \asymp \bn$. Isso tiraria $\lmaxx(\Gramhat)$ da
    Observação~\ref{rem:lenta-desenho}, mas não foi escrito.
\end{itemize}

\section{Conferência numérica}

\texttt{check/08-blocos.R} (imprime \texttt{OK}; cerca de $12$~min nesta máquina;
dependências \texttt{WaveBased} e \texttt{glmnet}; não chama
\texttt{wafc/R/}). Desenho de E1.4 a E1.6: Daublets com filtro $8$, $p = 2$
com $X_{1}\equiv 1$ e $X_{2} = 1/2 + \mathrm{Unif}(-1,1)$, $q = 2$, $\bU$
uniforme (logo $\gamma = 0{,}25$, $\kappa_{2} = 4/3$, $\Lambda = 1{,}458$),
$\sigma = 0{,}5$, $\alpha = 0{,}05$, $\bn = \lceil\log n\rceil$ e a forma
balanceada com pesos $\sqrt{|G|}$; o ajuste em blocos é o FISTA de
\texttt{check/08a-blocos.R} (KKT a menos de $10^{-7}$ de $\pen$ em todos os
ajustes) e o $\pen$ é o $\pen_{w}$ calibrado na forma fechada do
Lema~\ref{lem:calib-blocos}(ii). O que saiu em 2026-10-03:

\begin{itemize}[nosep, leftmargin=*]
  \item \textbf{Proposição~\ref{prop:balanceada}.} As três afirmações valem em
    $b = 2,\dots,200$ e $J = 1,\dots,14$. Com $J = 8$, o $\varrho^{2}$ dos pesos
    $\sqrt{|G|}$ é $1{,}667$, $1{,}571$, $1{,}571$ e $1{,}778$ nos $\bn$ de
    $n = 250$, $500$, $1000$ e $4000$.
  \item \textbf{Lema~\ref{lem:risco-pedacos}.} A cota vale nas $2\,280$
    sequências construídas ($19$ pares $(s,\pi)$ com $\pi$ até $\infty$,
    $b \in \{2,\dots,32\}$, $\eta$ de $10^{-9}$ a $10^{-1}$, $J = 20$, três
    formas que saturam a hipótese nível a nível), com razão máxima $0{,}823$:
    a cota está a um fator $1{,}2$ da pior sequência.
  \item \textbf{Os dois expoentes do Corolário~\ref{cor:adaptativa-blocos}},
    lidos no risco ideal dos pedaços finos com $60$ níveis: o expoente de
    $\eta/b$ é $2s/(2s+1)$ a menos de $10^{-3}$ nas duas formas e nos cinco
    pares testados; o de $b$ é exatamente $(2/\pi-1)/(2s+1)$ na espiga em
    $\pi < 2$ ($0{,}294$, $0{,}333$ e $0{,}128$), entre $-0{,}12$ e $0$ em
    $\pi \ge 2$, e $0$ na sequência espalhada no nível
    (Observação~\ref{rem:logaritmo}).
  \item \textbf{A aritmética.} A janela \eqref{eq:janela-blocos} é a do
    Corolário~5 em todo par da grade; dos pares testados, $(s,\pi) = (1;1)$,
    $(0{,}6;1{,}5)$, $(0{,}3;2)$, $(0{,}3;4)$ e $(0{,}3;\infty)$ estão no
    Corolário~\ref{cor:adaptativa-blocos} e fora do Teorema~2 de Klopp e
    Pensky.
  \item \textbf{O regime e a calibração} ($n = 250$ a $4000$ ao longo do
    $J_{n}$ do Teorema~\ref{thm:taxa-blocos} com $\seff = 1/2$, $40$ réplicas
    por $n$): a cobertura de $\Tw$ é $1{,}000$ com o $\pen_{w}$ fechado e com o
    $\lamG$ determinístico; o $\lamz{G}$ exato fica sob a forma fechada em
    toda réplica; $\smax^{2} \le 2B_{X}^{2}C_{U}$ sempre; $n(\lamGone)^{2}/\bn$
    fica entre $13{,}4$ e $17{,}8$. O evento $\ESig$ não vale nesses $n$
    ($\lmaxx(\Gramhat)$ mediano de $2{,}0$ a $2{,}7$); a $J = 4$ fixo,
    $\lmaxx(\Gramhat)$ desce de $1{,}85$ a $1{,}40$ e $\lmin(\Gramhat)$ sobe de
    $0{,}137$ a $0{,}234$ quando $n$ vai de $1000$ a $64000$, e só no último
    os dois estão no evento (Observação~\ref{rem:pre-blocos}).
  \item \textbf{Denso} ($\pi = 2$, $s = \seff = 1/2$, $n = 250$ a $8000$, $12$
    réplicas): as cotas do Teorema~\ref{thm:oraculo-blocos}(ii)--(iii) com
    $\bar\bbeta = \betastar$ e a do Corolário~\ref{cor:componentes-blocos}
    valem em todos os ajustes, com folga mínima de $3\,999$ vezes;
    $\gtil \ge \lmin(\Gramhat)$ sempre. A razão
    $\normn{\widehat f - f}^{2}/n^{-1/2}$ vai de $9{,}0$ a $10{,}6$ (máximo sobre
    mínimo $1{,}17$), e a inclinação de $\log(\text{erro})$ contra
    $\log n^{-1/2}$ é $1{,}027$, contra $1{,}194$ contra a forma com logaritmo
    $(\log n/n)^{1/2}$ do Teorema~2 de E1.6: o erro realizado segue a taxa sem
    logaritmo, o que E1.6 não conseguia separar. A soma dos erros das
    componentes fica a uma razão de $10{,}7$ a $18{,}7$ do alvo (máximo sobre
    mínimo $1{,}74$), crescendo até $n = 4000$ e estável depois.
  \item \textbf{Comprimível} ($\pi = 1$, $s = 1{,}2$, $J_{n}$ na ponta
    inferior da janela, $n = 250$ a $8000$, $10$ réplicas): o
    Teorema~\ref{thm:oraculo-blocos} vale com todos os comparadores truncados
    $\thetastar\mathbf{1}\{G \in T\}$, e as duas cotas do
    Corolário~\ref{cor:risco-ideal} e o Lema~\ref{lem:risco-pedacos} no
    $\thetastar$ do desenho valem em todos os ajustes. O $\gtil$ mediano vai
    de $0{,}012$ a $0{,}125$ e só chega a $\gamma/2$ em $n = 8000$; com ele, e
    com a constante conservadora da calibração ($n\eta_{n}/\bn$ perto de
    $200$, sob a cota $\kappa = 1017$ da prova), o $\Rid$ satura na energia
    total nesses $n$, e os expoentes ficam conferidos no item dos dois
    expoentes. O erro realizado fica a uma razão de $26{,}5$ a $44{,}6$ do alvo
    (máximo sobre mínimo $1{,}68$).
  \item \textbf{A variante branca} (Observação~\ref{rem:branca}; $n = 500$,
    $J = 5$, duas verdades, $40$ réplicas cada): $\operatorname{tr}(Q_{G}^{-1}
    \Psitil_{G})/|G| \le 0{,}9997$ e $\opnorm{Q_{G}^{-1/2}\Psitil_{G}
    Q_{G}^{-1/2}} \le 1$ em todos os pedaços; cobertura $1{,}000$ com o $\pen$
    pivotal; as três cotas com $q_{\max}\bar{\mathcal{W}}$ valem em todos os $80$
    ajustes, com folga mínima de $4\,899$ vezes; $q_{\max}/\lmaxx(\Gramhat) \le
    0{,}66$.
  \item \textbf{O limiarizado}: Parte~VI do script, relatada no adendo de
    E1.12 de \texttt{06-selecao-limiar.tex}.
  \item \textbf{Corolário~\ref{cor:lenta-blocos}} (Parte~VII, E1.13). A
    contagem \eqref{eq:contagem-lenta} fica acima do supremo exato de
    $\nGone{\thetastar}$ sobre a bola de Besov (nível a nível) em $12$ pares
    $(s,\pi)$ com $\pi$ até $\infty$, $b = 2,\dots,64$ e $J = 1,\dots,40$, com
    razão máxima $0{,}999$; num nível fino, o supremo cai como
    $b^{-\vartheta_{\pi}}$ à terceira casa. Os expoentes da cota, lidos no
    mínimo em $J$ ($2^{J} \le n/(\log n)^{2}$, $n$ de $10^{20}$ a $10^{300}$):
    o de $n$ é o do enunciado a $0{,}01$, nos blocos e no LASSO; o do
    logaritmo, nos blocos, fica entre $0{,}43$ e $0{,}52$ em~(i) (alvo $1/2$),
    entre $-0{,}002$ e $0{,}017$ em~(ii) com $\pi \ge 2$ (alvo $0$), e a menos
    de $0{,}05$ do alvo nos dois pares de~(ii) com $\pi < 2$, nos pesos~$1$ e
    $\sqrt{|G|}$; o da cota do LASSO (Proposição~4 de E1.6) fica a menos de
    $0{,}07$ do seu alvo, $1/2$ ou $2\seff/(1-2s+4\seff)$, e acima do dos blocos
    por pelo menos $0{,}28$ em~(ii). Numa grade de $156$ pares, $11$ ficam fora
    de $\seff > 1/4$ ou $\seff > s/2$, todos com $\pi < 2$, e neles a escolha
    de $J_{n}$ da Proposição~4 também não é admissível. Em ajustes ($\pen =
    \pen_{w}$ exato, $8$ réplicas por célula, os dois regimes e os dois
    pesos), em $\bU$ uniforme com $n = 250$ e $1000$, com $U_{1} \sim
    \mathrm{Unif}(0,1/2)$ ($n = 1000$, $J = 6$, colunas nulas) e com $d > n$
    ($n = 250$, $J = 7$, $d = 508$): cobertura de $\Tw$ igual a $1{,}000$,
    nenhuma violação do Teorema~\ref{thm:oraculo-blocos}(i) nem de
    \eqref{eq:lenta-nao-assint}, folga mínima de $27$ vezes, e $\gtil = 0$ (a
    menos de $10^{-14}$) nos dois últimos desenhos. No $\thetastar$ do
    desenho, $\norm{\thetastar}_{1}/\nGone{\thetastar}$ vai de $1{,}5$ a $1{,}9$
    em~(i) e de $2{,}6$ a $2{,}7$ em~(ii).
  \item \textbf{Corolário~\ref{cor:lenta-blocos}(iii) e (iv)} (Parte~VIII,
    E1.14). A cota com $\lamGp$, minimizada em $J$ com $2^{J} \le n^{2}$
    ($n$ de $10^{15}$ a $10^{140}$), em $17$ pares (quatro de (iii)(a), quatro
    de (iii)(b), cinco de (iv), quatro de (i) e (ii)): o expoente de $n$ é o do
    item a $0{,}01$ nos pesos~$1$ e $\sqrt{|G|}$; o do logaritmo, a $0{,}1$
    (por exemplo $0{,}21$ contra $0{,}20$ em $(s,\pi) = (0{,}4;1{,}25)$ e
    $0{,}95$ contra $1$ em $(0{,}5;2)$); o ganho sobre o LASSO no logaritmo é o
    da teoria em todo par; o minimizador cresce como o $J_{n}$ do enunciado
    (a $0{,}03$) e passa do teto exatamente nos pares de (iii) e de (iv) com
    $\seff < 1/4$; a cota no $J_{n}$ do enunciado tem o expoente da teoria; e
    com $J$ parado no teto e $\lamG$ o expoente é $-2\seff$, pior pela diferença
    da teoria. Numa grade de $183$ pares com $\seff > 0$, o expoente de $n$ é o
    do item do par, com desvio máximo de $0{,}002$. Acima do teto ($n = 250$,
    $J = 5$ a $9$, $10$ réplicas): $\smax^{2} \le 2B_{X}^{2}C_{U}\mu^{\Gcal}_{J}$
    e $\lmaxx(\Gramhat) \le \Lambda\mu^{\Gcal}_{J}$ em toda réplica (razões
    máximas $0{,}23$ e $0{,}035$; $c_{\mu} = 153$), $\pen_{w} \le \lamGp$ e
    cobertura de $\Tw$ igual a $1{,}000$ nos dois pesos; $\lmaxx(\Gramhat)$
    mediano vai de $2{,}4\Lambda$ a $14{,}7\Lambda$. Em ajustes ($n = 250$,
    $\pen = \pen_{w}$ exato, $4$ réplicas, os dois pesos) com $(s,\pi) =
    (0{,}6;1)$ em $J = 8$, $(0{,}4;1{,}25)$ e $(0{,}5;1{,}2)$ em $J = 7$, todos
    com $d > n$ e $\gtil = 0$: cobertura $1{,}000$, nenhuma violação do
    Teorema~\ref{thm:oraculo-blocos}(i) nem de \eqref{eq:lenta-nao-assint},
    folga mínima de $61$ vezes.
  \item \textbf{Corolário~\ref{cor:truncada-blocos}} (Parte~IX, E1.15). A
    contagem \eqref{eq:contagem-truncada-blocos} contra o supremo exato de
    $\mathcal{R}_{\Gcal,1}(\thetastar;t)$ na classe, na forma balanceada (nos
    níveis finos com $\pi \le 2$, uma espiga por pedaço reduz o problema ao
    coordenado com $k_{\lev}$ coordenadas; com $\pi > 2$, o espalhado contra
    Jensen; o pedaço grosso, exato), em $9$ pares com $s < 1/2$, $b = 2$ a
    $64$, $t$ de $10^{-1}$ a $10^{-40}$ e $J \le 400$: razão máxima à cota de
    $0{,}97$, e o supremo a pelo menos $0{,}79$ dela; cai como $t^{4s/(2s+1)}$ e
    como $b^{((2/\pi-1)_{+}-2s)/(2s+1)}$ (os dois a $0{,}002$). A cota com
    $\pen = (\mu^{\Gcal}_{J})^{1/2}\lamGone$ (pesos~$1$) minimizada em $J$ com
    $2^{J} \le n^{2}$ ($n$ de $10^{15}$ a $10^{140}$), em $16$ pares, ao lado da
    do LASSO truncado: o expoente de $n$ é o do corolário nos dois (a $0{,}01$);
    o do logaritmo, a $0{,}1$ em $n \ge 10^{60}$ (por exemplo $0{,}25$ contra
    $0{,}24$ em $(0{,}4;1{,}25)$ e $-0{,}01$ contra $0$ em $(0{,}3;2)$); o ganho
    no logaritmo sobre o LASSO truncado é o da teoria em todo par; o expoente é
    estritamente melhor que o do Corolário~\ref{cor:lenta-blocos} nos $8$
    pares com $s < 1/2$ e $\pi < 2$ e igual nos outros; nos $6$ pares de (iv)
    o minimizador passa do teto e segue o $J_{n}$ do enunciado (a $0{,}03$), e
    a cota no $J_{n}$ do enunciado fica a menos de um fator $7$ do mínimo. Numa
    grade de $212$ pares, desvio máximo de $0{,}005$ no expoente de $n$, nunca
    pior que o do Corolário~\ref{cor:lenta-blocos}. Em ajustes ($n = 250$, $4$
    réplicas, $\pen = \pen_{w}$ exato, pesos~$1$ e $\sqrt{|G|}$) com
    $(s,\pi) = (0{,}4;1{,}8)$ em $J = 5$, abaixo do teto, e $(0{,}4;1{,}25)$ e
    $(0{,}3;2)$ em $J = 7$, acima e com $\gtil = 0$: cobertura de $\Tw$ igual
    a $1{,}000$, nenhuma violação do Teorema~\ref{thm:oraculo-blocos}(i) no
    comparador truncado nem da cota de (i), folga mínima de $40$ vezes; a
    cadeia até a contagem vale no $\thetastar$ de cada ajuste; ficam no
    comparador $53\%$, $7\%$ e $5\%$ dos pedaços. A cota truncada é de
    $0{,}91$ a $0{,}94$ vezes a do Corolário~\ref{cor:lenta-blocos} acima do
    teto e de $1{,}28$ a $1{,}33$ vezes abaixo dele, pela mesma razão das
    constantes que no LASSO.
\end{itemize}

\section*{Referências}
\sloppy

Todas em \texttt{docs/referencias-verificadas.bib}.

\begin{itemize}[nosep, leftmargin=*]
  \item Bühlmann, P.\ e van de Geer, S. (2011). \emph{Statistics for
    High-Dimensional Data}. Springer. (\texttt{buhlmann2011statistics}.)
    Usada por meio de E1.5.
  \item Cai, T.~T. (1999). Adaptive wavelet estimation: a block thresholding
    and oracle inequality approach. \emph{The Annals of Statistics} 27(3),
    898--924. (\texttt{Cai-1999}, verificada em L5.) Teorema~4, eq.~(5.4),
    p.~908: a BlockJS está a $(\log n)^{(2/p-1)/(1+2\alpha)}$ do minimax em
    $\Besov{\alpha}{p}{q}$ com $1 \le p < 2$ e $\alpha \ge 1/p$, no modelo de
    sequência (conferido no PDF da editora em 2026-10-03).
  \item Hsu, D., Kakade, S.~M.\ e Zhang, T. (2012). A tail inequality for
    quadratic forms of subgaussian random vectors. \emph{Electronic
    Communications in Probability} 17, no.~52, 1--6.
    (\texttt{Hsu-Kakade-Zhang-2012}, verificada em L5.) Teorema~2.1, com
    $\mu = 0$ (Observação~2.2); o enunciado publicado toma
    $A \in \R^{n\times n}$ e $\Sigma = A\T A$ (conferido no PDF em
    2026-10-03).
  \item Klopp, O.\ e Pensky, M. (2015). Sparse high-dimensional varying
    coefficient model: nonasymptotic minimax study. \emph{The Annals of
    Statistics} 43(3), 1273--1299. (\texttt{Klopp-Pensky-2015}, verificada em
    L1, numeração conferida em L7.) Norma (3.1) e estimador (3.2)--(3.3),
    p.~1280; Teorema~1, eqs.~(3.6) e~(3.7), pp.~1281--1282; (3.13), (3.14),
    Teorema~2 e Observação~2, pp.~1283--1284; (3.19)--(3.21), p.~1285; (A1)(b)
    e (A3), pp.~1278--1279; (A6), p.~1278; o Lema~4 do suplemento
    (\texttt{10.1214/15-AOS1309SUPP}), eq.~(5.31), conferido no PDF em
    2026-10-01 e de novo em 2026-10-03.
  \item Lounici, K., Pontil, M., van de Geer, S.\ e Tsybakov, A.~B. (2011).
    Oracle inequalities and optimal inference under group sparsity.
    \emph{The Annals of Statistics} 39(4), 2164--2204.
    (\texttt{Lounici-Pontil-vandeGeer-Tsybakov-2011}, verificada em L3,
    numeração conferida em L7.) Lema~3.1 (p.~2173), Hipótese~3.1 e
    Teorema~3.1 (pp.~2176--2177), Teorema~7.1 (p.~2187).
\end{itemize}

\end{document}
